\documentclass[UTF8,12pt,a4paper,fontset=fandol]{ctexbook}


% 支持从项目根目录或当前笔记目录编译。
\makeatletter
\def\input@path{{../}}
\makeatother
\usepackage{researchnotes}


\title{复分析笔记}
\author{}
\date{\today}


\begin{document}


\maketitle
\tableofcontents
\clearpage


\chapter{绪论}
\label{chap:ca-introduction}

\section{研究对象与贯穿全书的问题}

\keyterm{复分析（complex analysis）}研究复变量函数的微分、积分、局部展开与大域性质。本书以单复变量为范围，从复平面的几何出发，逐步解释一个核心现象：一个函数只要在开集上处处复可微，就受到远强于实可微性的约束。它在每一点附近都能展开为收敛幂级数，其内部取值与边界积分相联系，其零点、奇点以及映射形状也具有严格的结构。

全书围绕四组问题展开：复导数为什么要求不同方向的差商给出同一个极限？曲线积分何时与路径无关，何时能够恢复区域内部的函数？局部级数如何控制零点和奇点，并帮助计算实积分？局部定义的函数怎样延拓、逼近，或者把一个区域映到另一个区域？这些问题依次连接分析、平面拓扑与几何，而留数计算只是其中一条应用线索。

\section{先修知识、课程范围与阅读层次}

默认读者掌握一元微积分、实数列与函数项级数、多元函数的偏导数和全微分，以及二维向量和线性映射。开集、紧集、连通性、一致收敛和曲线积分将在需要时复习；抽象拓扑、测度论和泛函分析不作为核心部分的先修课程。证明中使用紧性、极限交换或积分估计时，均需说明具体依据。

本书按本科复分析课程的核心内容组织，并设置进阶选读部分。第\ref{chap:ca-plane}章至第\ref{chap:ca-conformal}章构成首次学习的主体；第\ref{chap:ca-normal}章至第\ref{chap:ca-approximation}章进一步讨论存在性、延拓和构造问题。黎曼曲面只作入门，椭圆函数、模形式、多复变量与完整的值分布理论不在本书正文范围内。

本书将基本概念、核心证明与典型计算结合组织。正文优先说明结论为什么成立、条件在哪里使用；技术性推广另作提示。第一篇由复平面逐步进入积分理论，后续各篇仍保留章纲，依次展开函数结构、共形几何与大域理论。

\section{课程总纲}

下表给出全书的主题、核心问题和学习层次。“进阶”部分适合完成核心课程后选读。

\begin{longtable}{@{}p{0.07\textwidth}p{0.28\textwidth}p{0.49\textwidth}p{0.08\textwidth}@{}}
\toprule
章节 & 主题 & 核心问题 & 层次 \\
\midrule
\endhead
\ref{chap:ca-plane} & 复数与复平面 & 复数运算的几何意义；区域与曲线的基本语言 & 核心 \\
\ref{chap:ca-holomorphic} & 复可微性与全纯函数 & 复导数与实全微分的联系；柯西--黎曼条件 & 核心 \\
\ref{chap:ca-elementary} & 初等函数与分支 & 指数、对数、幂函数；多值性从何而来 & 核心 \\
\ref{chap:ca-integration} & 复积分与原函数 & 参数化、积分估计、路径无关与绕数 & 核心 \\
\ref{chap:ca-cauchy} & 柯西理论 & 积分定理、积分公式及其适用的区域条件 & 核心 \\
\ref{chap:ca-series} & 幂级数与解析性 & 全纯性为何等价于局部幂级数表示 & 核心 \\
\ref{chap:ca-local} & 零点与局部结构 & 恒等定理、开映射、最大模与增长约束 & 核心 \\
\ref{chap:ca-singularities} & 洛朗级数与奇点 & 环域展开、奇点分类与无穷远点 & 核心 \\
\ref{chap:ca-residues} & 留数及积分应用 & 奇点如何决定闭路积分；实积分如何处理 & 核心 \\
\ref{chap:ca-argument} & 辐角原理与零点计数 & 不求出零点也能计数；扰动如何保持零点数 & 核心 \\
\ref{chap:ca-harmonic} & 调和函数 & 调和共轭、泊松公式与圆盘边值问题 & 核心 \\
\ref{chap:ca-conformal} & 共形映射 & 局部保角与整体双全纯；典型区域的变换 & 核心 \\
\ref{chap:ca-normal} & 正规族与黎曼映射 & 函数列的紧性如何产生共形映射 & 进阶 \\
\ref{chap:ca-continuation} & 解析延拓与黎曼曲面 & 局部函数如何延续；沿路径绕行会发生什么 & 进阶 \\
\ref{chap:ca-entire} & 整函数与亚纯函数构造 & 怎样指定零点或主部并构造函数 & 进阶 \\
\ref{chap:ca-approximation} & 全纯逼近与综合专题 & 怎样用有理函数逼近；如何综合使用全书工具 & 进阶 \\
\bottomrule
\end{longtable}

\section{知识依赖与学习路线}

基础主线是复数与拓扑语言、复可微性、曲线积分、柯西理论、幂级数、零点与奇点，最后进入留数和辐角原理。初等函数在主线前段提供可反复使用的例子；调和函数和共形映射在主线后段汇合几何与分析。

为避免证明上的循环依赖，第\ref{chap:ca-holomorphic}章先以复可微定义全纯性；第\ref{chap:ca-elementary}章借助实指数和实三角函数构造复指数；第\ref{chap:ca-cauchy}章从复可微性证明柯西积分理论；到第\ref{chap:ca-series}章再证明全纯函数具有局部幂级数展开。不能提前用“全纯函数已经可以展开为幂级数”证明柯西定理。

初次学习可顺序完成核心各章，并在每章检查三个层面：能否准确复述条件，能否解释证明中的关键步骤，能否处理一个不满足条件的反例。完成核心部分后，进阶各章按正文顺序学习；若主要关心函数构造，也可在第\ref{chap:ca-continuation}章之后直接集中学习第\ref{chap:ca-entire}章。

\section{记号约定}

以 $\mathbb{R}$、$\mathbb{C}$ 分别表示实数域、复数域，$i$ 表示满足 $i^2=-1$ 的虚数单位。通常写 $z=x+iy$，其中 $x,y\in\mathbb{R}$；$\overline z$、$|z|$、$\operatorname{Re}z$、$\operatorname{Im}z$ 分别表示共轭、模、实部和虚部。复函数通常写成 $f=u+iv$，其中 $u,v$ 是实值函数。

区域 $\Omega$ 指非空连通开集，$D(a,r)=\{z\in\mathbb{C}:|z-a|<r\}$ 表示以 $a\in\mathbb{C}$ 为中心、半径 $r>0$ 的开圆盘，$\mathbb{D}=D(0,1)$。$\partial\Omega$ 表示边界，$\overline\Omega$ 表示闭包，$K\Subset\Omega$ 表示 $K$ 是包含于 $\Omega$ 的紧集。曲线积分默认沿分段连续可微曲线进行；正向边界按行进方向使区域位于左侧，因此多连通区域的内边界取顺时针方向。

章末习题涵盖基本计算、证明与反例，解答或提示列于附录\ref{chap:ca-solutions}。符号 $\log r$（$r>0$）表示实自然对数；复对数的分支另行定义。未经特别说明，函数等式均在两侧有定义的共同区域内讨论。

\section{参考学习材料}

课程覆盖范围可对照 MIT 的\href{https://ocw.mit.edu/courses/18-04-complex-variables-with-applications-spring-2018/pages/lecture-notes/}{18.04：Complex Variables with Applications（2018）}与\href{https://ocw.mit.edu/courses/18-112-functions-of-a-complex-variable-fall-2008/pages/lecture-notes/}{18.112：Functions of a Complex Variable（2008）}官方讲义目录。前者适合配合计算与应用，后者可用于核对积分理论、无穷乘积、正规族和黎曼映射等主题。进阶方向可参照 Harvard 的\href{https://people.math.harvard.edu/~ctm/home/text/class/harvard/213a/06/html/syl.html}{Math 213a（2006）课程大纲}。本书的章节顺序按上述知识依赖独立组织，不将这些课程的全部专题纳入正文。


\part{基本对象与积分理论}

\chapter{复数与复平面}
\label{chap:ca-plane}

实数轴上的微积分研究一维变化，复分析则把二维平面与一种特殊的乘法结合起来。乘法决定了复导数的几何性质，平面的拓扑决定了积分与函数分支的整体行为。本章建立这两方面的基本语言。

\section{复数的代数结构}

\begin{definition}[复数]
在 $\mathbb{R}^2$ 上定义
\[
 (x,y)+(u,v)=(x+u,y+v),\qquad
 (x,y)(u,v)=(xu-yv,xv+yu).
\]
所得数系称为\keyterm{复数域（complex field）}，记为 $\mathbb{C}$。把 $(x,0)$ 与实数 $x$ 等同，记 $i=(0,1)$，则 $i^2=-1$，且 $(x,y)=x+iy$。
\end{definition}

加法、乘法的结合律与分配律可按上述定义展开验证；零元为 $(0,0)$，单位元为 $(1,0)$。若 $z=x+iy\ne0$，则
\[
 z^{-1}=\frac{x-iy}{x^2+y^2}.
\]
因此每个非零复数都有乘法逆元。对 $z=x+iy$，称 $x=\operatorname{Re}z$、$y=\operatorname{Im}z$ 为实部、虚部，定义\keyterm{共轭（complex conjugate）}$\overline z=x-iy$ 与\keyterm{模（modulus）}$|z|=\sqrt{x^2+y^2}$。直接计算得
\begin{equation}
 z\overline z=|z|^2,\quad
 \overline{zw}=\overline z\,\overline w,\quad
 |zw|=|z||w|,\quad
 \frac{z}{w}=\frac{z\overline w}{|w|^2}\quad(w\ne0).
 \label{eq:ca-modulus-rules}
\end{equation}
其中模的乘法公式由 $|zw|^2=(z\overline z)(w\overline w)$ 及模非负得到。

\begin{proposition}[模的不等式]
\label{prop:ca-triangle}
对任意 $z,w\in\mathbb{C}$，有
\[
 |z+w|\le|z|+|w|,\qquad
 \bigl||z|-|w|\bigr|\le|z-w|.
\]
第一个不等式取等号，当且仅当 $z,w$ 中至少一个为零，或二者均非零且 $z/w$ 为正实数。
\end{proposition}
\begin{proof}
由 $\operatorname{Re}(z\overline w)\le|z\overline w|=|z||w|$，
\[
 |z+w|^2=|z|^2+2\operatorname{Re}(z\overline w)+|w|^2
 \le(|z|+|w|)^2.
\]
取非负平方根即得第一式。二者非零时，等号等价于 $z\overline w$ 是正实数，也等价于 $z/w>0$。再由 $|z|\le|z-w|+|w|$ 及交换 $z,w$ 后的同一不等式得到第二式。
\end{proof}

复数的模可以比较大小，复数本身却不存在与域运算相容的全序：在有序域中非零元素的平方均为正，因而应有 $1>0$ 与 $i^2=-1>0$，矛盾。表达式 $z>0$ 仅在已经说明 $z$ 为实数时使用。

\section{极形式、辐角与开方}

非零复数 $z=x+iy$ 可写成
\[
 z=r(\cos\theta+i\sin\theta),\qquad r=|z|>0.
\]
满足此式的 $\theta$ 称为 $z$ 的\keyterm{辐角（argument）}。若 $\theta_0$ 是一个辐角，则全部辐角为 $\theta_0+2k\pi$（$k\in\mathbb{Z}$），记作集合 $\arg z$。区间 $(-\pi,\pi]$ 内的唯一辐角称为\keyterm{主辐角（principal argument）}，记作 $\operatorname{Arg}z$。零没有辐角；主辐角也不能在整个去心平面上连续。

由三角函数加法公式，两个非零复数相乘时，模相乘而辐角相加。因此乘以固定复数 $r(\cos\theta+i\sin\theta)$，在平面上表示以原点为中心旋转 $\theta$ 后放大 $r$ 倍。归纳并结合求逆可得棣莫弗公式（de Moivre's formula）：对每个整数 $n$，
\begin{equation}
 (\cos\theta+i\sin\theta)^n=\cos(n\theta)+i\sin(n\theta).
 \label{eq:ca-demoivre}
\end{equation}

\begin{proposition}[复数开方]
设 $n\ge1$ 为整数，$w=\rho(\cos\phi+i\sin\phi)\ne0$。方程 $z^n=w$ 恰有 $n$ 个不同的解：
\begin{equation}
 z_k=\rho^{1/n}\left(\cos\frac{\phi+2k\pi}{n}
       +i\sin\frac{\phi+2k\pi}{n}\right),\quad 0\le k<n.
 \label{eq:ca-roots}
\end{equation}
\end{proposition}
\begin{proof}
写 $z=r(\cos\theta+i\sin\theta)$。由式\eqref{eq:ca-demoivre}，原方程等价于 $r^n=\rho$ 及 $n\theta-\phi\in2\pi\mathbb{Z}$。模取唯一正根，辐角模 $2\pi$ 恰给出上述 $n$ 个值。
\end{proof}

\begin{example}[单位根之和]
设 $n\ge2$，$\omega=\cos(2\pi/n)+i\sin(2\pi/n)$。全部 $n$ 次单位根为 $1,\omega,\ldots,\omega^{n-1}$，它们位于单位圆上等分圆周。由于 $\omega\ne1$，
\[
 1+\omega+\cdots+\omega^{n-1}=\frac{1-\omega^n}{1-\omega}=0.
\]
几何上，这说明正 $n$ 边形各顶点的位置向量之和为零。
\end{example}

\section{复平面的拓扑基础}

模给出欧氏距离 $d(z,w)=|z-w|$，因此 $\mathbb{C}$ 的距离与 $\mathbb{R}^2$ 完全相同。以下概念均相对于这一距离。

\begin{definition}[开集、闭集与边界]
若每个 $a\in U$ 都有某个圆盘 $D(a,r)\subset U$，则称 $U$ 为\keyterm{开集（open set）}；补集为开集的集合称为\keyterm{闭集（closed set）}。若 $a$ 的每个圆盘都与集合 $E$ 相交，则 $a$ 属于 $E$ 的闭包 $\overline E$；若每个圆盘同时与 $E$ 及其补集相交，则 $a$ 属于边界 $\partial E$。若每个去心圆盘都与 $E$ 相交，则称 $a$ 为 $E$ 的\keyterm{聚点（accumulation point）}。
\end{definition}

闭包允许用集合中的点本身检验，聚点则要求任意小邻域内还有别的点。例如单点集是闭集，但没有聚点。数列 $z_n=x_n+iy_n$ 收敛到 $z=x+iy$，当且仅当 $x_n\to x$ 且 $y_n\to y$，因为
\[
 |x_n-x|,\ |y_n-y|\le|z_n-z|\le|x_n-x|+|y_n-y|.
\]
同样的估计说明复数列是柯西列当且仅当两个坐标列都是柯西列。由实数完备性，$\mathbb{C}$ 也是完备的。

集合 $K$ 称为\keyterm{紧集（compact set）}，若它的每个开覆盖都有有限子覆盖。实欧氏空间中的海涅--博雷尔定理（Heine--Borel theorem）直接给出：复平面中的集合紧，当且仅当它有界且闭。例如闭圆盘紧，开圆盘不紧。若 $K\ne\varnothing$ 紧、$U$ 开且 $K\subset U$，则存在 $\delta>0$，使每个 $z\in K$ 都有 $D(z,\delta)\subset U$。证明是为每个点选取一个圆盘，先将半径减半，再对这些半径减半的圆盘取有限子覆盖。

\begin{definition}[连通与区域]
集合 $E$ 称为\keyterm{连通（connected）}，若不能写成两个不交、非空、在 $E$ 中相对开集合的并；这里“在 $E$ 中相对开”指形如 $E\cap U$，其中 $U$ 在 $\mathbb{C}$ 中开。非空连通开集称为\keyterm{区域（domain）}。若 $E$ 中任意两点可由留在 $E$ 内的连续曲线连接，则称 $E$ 道路连通（path connected）。
\end{definition}

\begin{proposition}
\label{prop:ca-polygonal-connected}
复平面中每个区域内的任意两点均可用有限折线连接。
\end{proposition}
\begin{proof}
固定 $a\in\Omega$，令 $A$ 为可由 $\Omega$ 内有限折线与 $a$ 相连的点集。对任意 $z\in\Omega$，取 $D(z,r)\subset\Omega$。若 $z\in A$，把一小段直线接在已有折线之后，可得 $D(z,r)\subset A$；若 $z\notin A$，该圆盘也不能含有 $A$ 中的点，否则可再接线段到达 $z$。故 $A$ 及 $\Omega\setminus A$ 在 $\Omega$ 中都开。由连通性及 $a\in A$，只能有 $A=\Omega$。
\end{proof}

道路连通集合必连通，因为区间的连续像连通。结合命题\ref{prop:ca-polygonal-connected}，开集的连通性与道路连通性等价；这不是任意集合都具有的性质。

\section{曲线、同伦与单连通性}

一条\keyterm{曲线（curve）}是连续映射 $\gamma:[0,1]\to\mathbb{C}$。若 $\gamma(0)=\gamma(1)$，称其闭合；闭曲线若除了这对端点外不重复取值，称为简单闭曲线（simple closed curve）。积分中还要求曲线分段 $C^1$，即把参数区间有限分段后，各段上 $\gamma$ 连续可微。曲线的像是点集，曲线本身还包含行进方向和重复经过的次数，二者不能混淆。

\begin{definition}[同伦与单连通]
同端点曲线 $\gamma_0,\gamma_1:[0,1]\to\Omega$ 的\keyterm{定端点同伦（homotopy relative to endpoints）}是连续映射 $H:[0,1]^2\to\Omega$，满足
\[
 H(s,0)=\gamma_0(s),\quad H(s,1)=\gamma_1(s),\quad
 H(0,t)=\gamma_0(0),\quad H(1,t)=\gamma_0(1).
\]
区域 $\Omega$ 称为\keyterm{单连通（simply connected）}，若每条闭曲线都能在 $\Omega$ 内连续收缩为常值曲线，即存在连续 $H$，使 $H(s,0)=\gamma(s)$、$H(s,1)=a$ 且 $H(0,t)=H(1,t)$。这里中间各曲线均闭合，但公共端点可以随 $t$ 移动。
\end{definition}

若区域关于 $a$ 星形，即对每个 $z\in\Omega$，线段 $[a,z]\subset\Omega$，则 $H(s,t)=(1-t)\gamma(s)+ta$ 给出所需收缩。因此星形区域单连通，特别地，圆盘、半平面和整个复平面均单连通。圆环 $r<|z|<R$（$0<r<R$）与去心平面则不单连通：绕孔一周的曲线不能在区域内缩成一点。其严格证明将由绕数的同伦不变性给出，见第\ref{chap:ca-integration}章和第\ref{chap:ca-cauchy}章。

\begin{example}[几个典型区域]
圆盘、圆环和裂平面 $\mathbb{C}\setminus(-\infty,0]$ 都开且连通。圆环不单连通；裂平面关于点 $1$ 星形，故单连通。两个不交开圆盘的并是开集，却不连通，因此不是区域。“有界”“连通”和“单连通”是不同的要求。
\end{example}

\section{扩充复平面与球面表示}

为统一研究有限点与无穷远处的行为，在 $\mathbb{C}$ 中添加一个点 $\infty$，得到\keyterm{扩充复平面（extended complex plane）}$\widehat{\mathbb{C}}$。无穷远点的基本邻域定义为 $\{z:|z|>R\}\cup\{\infty\}$，其中 $R>0$；于是 $z_n\to\infty$ 当且仅当 $|z_n|\to\infty$，而不要求辐角趋向某一方向。

令 $S^2=\{(X,Y,Z)\in\mathbb{R}^3:X^2+Y^2+Z^2=1\}$，北极为 $N=(0,0,1)$。球极投影（stereographic projection）将 $S^2\setminus\{N\}$ 映到复平面：
\begin{equation}
 z=\frac{X+iY}{1-Z},\qquad
 (X,Y,Z)=\frac{1}{1+|z|^2}
 \bigl(2\operatorname{Re}z,2\operatorname{Im}z,|z|^2-1\bigr).
 \label{eq:ca-stereographic}
\end{equation}
这两个公式互为逆映射；当 $|z|\to\infty$ 时，右侧趋于 $N$。把 $N$ 对应于 $\infty$，就得到\keyterm{黎曼球面（Riemann sphere）}的表示。它也是紧空间：含有 $\infty$ 的一个开集覆盖某个大圆盘的外部，剩余闭圆盘可用有限个开集覆盖。

无穷远点附近的局部坐标取为 $w=1/z$，并规定 $w(\infty)=0$。因而研究 $f$ 在无穷远处的行为，可以转为研究 $f(1/w)$ 在 $w=0$ 附近的行为。$\infty$ 是新增的几何点，不是一个可任意参与域运算的复数；例如 $\infty-\infty$ 没有由此定义。

\section{例题与习题}

\begin{example}[阿波罗尼斯圆]
设 $a,b\in\mathbb{C}$、$a\ne b$，$k>0$。求满足 $|z-a|=k|z-b|$ 的点集。
\end{example}
\begin{solve}
平方并展开，得到
\[
 (1-k^2)|z|^2-2\operatorname{Re}\bigl(z(\overline a-k^2\overline b)\bigr)
 +|a|^2-k^2|b|^2=0.
\]
若 $k=1$，这表示到 $a,b$ 等距的垂直平分线。若 $k\ne1$，配方得圆
\[
 \left|z-\frac{a-k^2b}{1-k^2}\right|=\frac{k|a-b|}{|1-k^2|}.
\]
因此距离之比固定的轨迹可以是圆，也可以在特殊参数下退化为直线。
\end{solve}

\begin{enumerate}
  \item\label{ex:ca-roots-sum} 求 $z^4=-4$ 的全部解。设 $\omega$ 为本章例题中的 $n$ 次单位根，证明对任意整数 $m$，$\sum_{k=0}^{n-1}\omega^{mk}$ 等于 $n$ 或 $0$，并给出两种情形的条件。
  \item\label{ex:ca-parallelogram} 证明 $|z+w|^2+|z-w|^2=2(|z|^2+|w|^2)$，并解释其几何意义。
  \item\label{ex:ca-domains} 分别判断 $\{z:|z|<1\}$、$\{z:0<|z|<1\}$、$\mathbb{C}\setminus\mathbb{R}$ 和 $\{z:\operatorname{Re}z>0\}$ 的连通性与单连通性。
  \item\label{ex:ca-separation} 设非空紧集 $K$ 包含于开集 $U\ne\mathbb{C}$，证明 $\inf\{|z-w|:z\in K,\ w\notin U\}>0$。
\end{enumerate}

\chapter{复可微性与全纯函数}
\label{chap:ca-holomorphic}

在实变量中，导数描述沿一条直线的局部变化；在复变量中，增量可从任意方向趋于零，而差商必须趋于同一个复数。这一要求使一般二维实线性映射缩小为复数乘法，是全纯函数刚性的起点。

\section{复函数的极限与连续性}

设 $f:E\to\mathbb{C}$，$z_0$ 为 $E$ 的聚点。称 $\lim_{z\to z_0}f(z)=L$，若对每个 $\varepsilon>0$，存在 $\delta>0$，使 $z\in E$ 且 $0<|z-z_0|<\delta$ 时都有 $|f(z)-L|<\varepsilon$。若 $z_0\in E$，把条件改为 $|z-z_0|<\delta$ 且取 $L=f(z_0)$，就得到在 $z_0$ 处连续的定义。

写 $f(x+iy)=u(x,y)+iv(x,y)$，则复极限存在当且仅当两个实分量的二元极限都存在。由此，极限的四则运算法则、连续映射保持紧性，以及紧集上的连续函数一致连续等结论都可从实分析得到。特别地，若 $f$ 在非空紧集 $K$ 上连续，则 $|f|$ 在 $K$ 上取得最大值与最小值。

\begin{example}[路径检验的限度]
令 $g(x+iy)=x^2y/(x^4+y^2)$（$(x,y)\ne(0,0)$）。沿任意过原点的直线，其极限均为零：当 $y=mx$ 且 $m\ne0$ 时，$g=mx/(x^2+m^2)\to0$，水平线与竖直线上也为零；但沿 $y=x^2$，$g=1/2$。所以即使检验了所有直线，仍不足以证明二元极限存在。
\end{example}

\section{复导数与全纯性的定义}

\begin{definition}[复导数与全纯函数]
设 $f$ 定义于 $z_0$ 的某个开邻域。若极限
\begin{equation}
  f'(z_0)=\lim_{h\to0}\frac{f(z_0+h)-f(z_0)}{h},\qquad h\in\mathbb{C}\setminus\{0\}
  \label{eq:ca-complex-derivative}
\end{equation}
存在，则称 $f$ 在 $z_0$ \keyterm{复可微（complex differentiable）}。若 $f$ 在开集 $U$ 的每一点复可微，称 $f$ 在 $U$ 上\keyterm{全纯（holomorphic）}；在整个 $\mathbb{C}$ 上全纯的函数称为\keyterm{整函数（entire function）}。
\end{definition}

复可微等价于一阶展开
\begin{equation}
 f(z_0+h)=f(z_0)+f'(z_0)h+o(|h|),\qquad h\to0,
 \label{eq:ca-first-order}
\end{equation}
其中 $o(|h|)$ 表示除以 $|h|$ 后趋于零的复数余项。此式立刻推出连续性，也给出求导法则：在相应导数存在且分母不为零处，
\[
 (f+g)'=f'+g',\quad (fg)'=f'g+fg',\quad
 (f/g)'=\frac{f'g-fg'}{g^2},\quad
 (g\circ f)'(z)=g'(f(z))f'(z).
\]
例如乘积法则由两份一阶展开相乘得到；链式法则则先令 $k=f(z_0+h)-f(z_0)=O(|h|)$，再把 $g(f(z_0)+k)$ 展开。于是多项式为整函数，有理函数在分母非零处全纯。

\begin{example}[点态可微与全纯]
函数 $f(z)=\overline z$ 的差商为 $\overline h/h$，沿实轴与虚轴分别等于 $1$、$-1$，故处处不可复微。对 $q(z)=|z|^2$，在 $z_0$ 处的差商为
\[
 \overline{z_0}+z_0\frac{\overline h}{h}+\overline h.
\]
两条坐标轴上的极限相等当且仅当 $z_0=0$；此时差商趋于零。所以 $q$ 仅在原点复可微，不在任何非空开集上全纯。
\end{example}

\section{柯西--黎曼方程及充分条件}

\begin{theorem}[复可微的实变量判别]
\label{thm:ca-cr}
设 $f=u+iv$ 定义于 $z_0=x_0+iy_0$ 的邻域。$f$ 在 $z_0$ 复可微，当且仅当 $(u,v)$ 在 $(x_0,y_0)$ 实可微，且满足\keyterm{柯西--黎曼方程（Cauchy--Riemann equations）}
\begin{equation}
 u_x=v_y,\qquad u_y=-v_x.
 \label{eq:ca-cr}
\end{equation}
此时 $f'(z_0)=u_x+iv_x=v_y-iu_y$，所有偏导数均在 $(x_0,y_0)$ 取值。
\end{theorem}
\begin{proof}
若复导数为 $A=a+ib$，式\eqref{eq:ca-first-order} 给出的实线性主部是
\[
 A(\Delta x+i\Delta y)=(a\Delta x-b\Delta y)
       +i(b\Delta x+a\Delta y),
\]
因此实微分存在，且 $u_x=v_y=a$、$v_x=-u_y=b$。反之，实可微性给出
\[
 f(z_0+h)-f(z_0)=(u_x+iv_x)\Delta x+(u_y+iv_y)\Delta y+o(|h|).
\]
代入式\eqref{eq:ca-cr}，右侧等于 $(u_x+iv_x)h+o(|h|)$，除以 $h$ 即得结论。
\end{proof}

特别地，若 $u,v$ 在开集 $U$ 上有连续的一阶偏导数，并在每点满足式\eqref{eq:ca-cr}，则 $f$ 在 $U$ 上全纯。连续偏导数用于保证实可微性；若已知实可微，则不必另加偏导数连续。

\begin{example}[仅有偏导数的反例]
定义
\[
 f(x+iy)=\begin{cases}
 \dfrac{xy}{x^2+y^2}(x+iy),&(x,y)\ne(0,0),\\
 0,&(x,y)=(0,0).
 \end{cases}
\]
因 $|xy|\le(x^2+y^2)/2$，有 $|f(z)|\le|z|/2$，故它在原点连续。两个坐标轴上函数恒为零，因此实部、虚部的四个偏导数在原点均为零，满足柯西--黎曼方程；但 $f(h)/h$ 沿实轴等于零，沿 $h=t(1+i)$ 等于 $1/2$。缺少的是实可微性，而非连续性。
\end{example}

\section{复导数的线性与几何解释}

复可微函数的雅可比矩阵为
\[
 Df(z_0)=\begin{pmatrix}a&-b\\ b&a\end{pmatrix},\qquad
 f'(z_0)=a+ib.
\]
若 $f'(z_0)\ne0$，这一主部把每个方向旋转同一角度、伸缩同一倍数，且 $\det Df(z_0)=|f'(z_0)|^2>0$。局部保角的严格结论将在第\ref{chap:ca-conformal}章给出；这里的线性近似尚不能保证映射在整个定义域上一一对应。

对实可微的复值函数，可定义\keyterm{维尔廷格导数（Wirtinger derivatives）}
\[
 \frac{\partial f}{\partial z}=\frac12(f_x-if_y),\qquad
 \frac{\partial f}{\partial\overline z}=\frac12(f_x+if_y).
\]
实微分可写为 $df=(\partial f/\partial z)\,dz+(\partial f/\partial\overline z)\,d\overline z$。定理\ref{thm:ca-cr} 因而等价于：实可微时，复可微恰好要求 $\partial f/\partial\overline z=0$，且复导数为 $\partial f/\partial z$。这是一种简洁记号，不能用来省略实可微假设。

\begin{proposition}
\label{prop:ca-zero-derivative}
若 $f$ 在区域 $\Omega$ 上全纯且 $f'\equiv0$，则 $f$ 为常数。特别地，区域上的实值全纯函数为常数。
\end{proposition}
\begin{proof}
由柯西--黎曼方程，$f'=0$ 使 $u,v$ 的四个一阶偏导数全为零。在区域内任一线段上，把 $u,v$ 与线段参数复合，实链式法则与一元中值定理说明二者不变。由命题\ref{prop:ca-polygonal-connected}，任意两点间可用有限折线连接，故 $f$ 为常数。若 $f$ 实值，则 $v\equiv0$，柯西--黎曼方程直接给出 $u_x=u_y=0$。
\end{proof}

\section{解析性与调和性的预告}

若 $f$ 在每点 $a$ 附近可写为收敛幂级数 $\sum_{n=0}^{\infty}c_n(z-a)^n$，则称它\keyterm{解析（analytic）}。单复变量中，解析性与全纯性等价；这不是定义的同义改写，而是第\ref{chap:ca-cauchy}章与第\ref{chap:ca-series}章要证明的核心结论。当前只由一次复可微，不能直接把无穷阶导数或幂级数展开作为已经证明的事实。

若全纯函数 $f=u+iv$ 的实部、虚部另外具有连续二阶偏导数，则由混合偏导可交换，
\[
 u_{xx}+u_{yy}=v_{yx}-v_{xy}=0,\qquad
 v_{xx}+v_{yy}=-u_{yx}+u_{xy}=0.
\]
满足拉普拉斯方程 $\Delta u=u_{xx}+u_{yy}=0$ 的二阶连续可微实函数称为\keyterm{调和函数（harmonic function）}。柯西积分公式将保证所有全纯函数都具有所需光滑性，从而这一额外假设最终可去掉。

\section{例题与习题}

\begin{example}[识别一个整函数]
设 $u=x^3-3xy^2$、$v=3x^2y-y^3$。求 $f=u+iv$ 的导数。
\end{example}
\begin{solve}
四个偏导数连续，且 $u_x=v_y=3x^2-3y^2$、$u_y=-v_x=-6xy$，故 $f$ 整，$f'=3(x^2-y^2)+6ixy=3z^2$。直接展开 $(x+iy)^3$ 也可识别 $f=z^3$；两种方法分别体现实变量判别和复代数计算。
\end{solve}

\begin{enumerate}
  \item\label{ex:ca-affine} 设 $f(z)=az+b\overline z+c$，其中 $a,b,c\in\mathbb{C}$。求其全部复可微点，并说明何时为整函数。
  \item\label{ex:ca-c1-criterion} 求实数 $\alpha,\beta$，使 $f(x+iy)=x^2+\alpha y^2+i\beta xy$ 为整函数，并求 $f'$。
  \item\label{ex:ca-real-constant} 设 $f$ 全纯，定义域为 $\{z:|\operatorname{Re}z|>1\}$，且 $f$ 实值。它是否必须在整个定义域上取同一常数？
  \item\label{ex:ca-modulus-constant} 若区域上的全纯函数满足 $|f|\equiv c$（$c\ge0$），利用柯西--黎曼方程证明 $f$ 为常数，不使用最大模原理。
\end{enumerate}

\chapter{初等全纯函数与分支}
\label{chap:ca-elementary}

多项式和有理函数由代数运算产生；指数、对数和非整数幂则把复分析中的周期与分支问题呈现出来。本章用实指数和实三角函数构造复指数，不预先使用复幂级数理论。

\section{复指数函数}

\begin{definition}[复指数]
对 $z=x+iy$，定义\keyterm{复指数函数（complex exponential）}
\begin{equation}
 e^z=\exp z=e^x(\cos y+i\sin y).
 \label{eq:ca-exponential}
\end{equation}
\end{definition}
其实部、虚部具有连续偏导数且满足柯西--黎曼方程，故 $e^z$ 是整函数，并有 $(e^z)'=e^z$。由实指数与三角函数的加法公式，
\[
 e^{z+w}=e^ze^w,\qquad |e^z|=e^{\operatorname{Re}z},\qquad e^z\ne0.
\]
当 $x=0$ 时得到欧拉公式（Euler's formula）$e^{iy}=\cos y+i\sin y$，从此极形式可简记为 $z=re^{i\theta}$。

\begin{proposition}
\label{prop:ca-exp-fibres}
$e^{z_1}=e^{z_2}$ 当且仅当 $z_1-z_2\in2\pi i\mathbb{Z}$。对每个 $w\ne0$，方程 $e^z=w$ 的全部解为
\[
 z=\log|w|+i(\theta+2k\pi),\qquad k\in\mathbb{Z},\quad \theta\in\arg w.
\]
\end{proposition}
\begin{proof}
由模相等知 $\operatorname{Re}z_1=\operatorname{Re}z_2$，再由单位圆上两点相等，虚部相差 $2\pi$ 的整数倍。第二个结论由比较式\eqref{eq:ca-exponential} 的模与辐角得到。
\end{proof}
所以指数的像集是 $\mathbb{C}\setminus\{0\}$，在整个平面上不单射。竖直线 $x=c$ 映为半径 $e^c$ 的圆，水平线 $y=c$ 映为辐角为 $c$ 的射线；竖直方向相隔 $2\pi$ 的点有同一个像。

\section{三角函数与双曲函数}

定义
\[
 \sin z=\frac{e^{iz}-e^{-iz}}{2i},\quad
 \cos z=\frac{e^{iz}+e^{-iz}}2,\quad
 \sinh z=\frac{e^z-e^{-z}}2,\quad
 \cosh z=\frac{e^z+e^{-z}}2.
\]
这四个函数均为整函数，在实轴上与相应的实函数一致。对指数求导或直接相乘，得到
\[
 (\sin z)'=\cos z,\quad(\cos z)'=-\sin z,\quad
 \sin^2z+\cos^2z=1,\quad \cosh^2z-\sinh^2z=1.
\]
加法公式同样由指数乘法公式推出。由 $\sin z=0\iff e^{2iz}=1$ 及命题\ref{prop:ca-exp-fibres}，正弦的全部零点是 $k\pi$；余弦的全部零点是 $\pi/2+k\pi$，其中 $k\in\mathbb{Z}$。这些零点处的一阶导数均不为零。

三角函数在复平面上不再有实轴上的有界性。例如 $\sin(iy)=i\sinh y$、$\cos(iy)=\cosh y$，其模可随 $|y|\to\infty$ 而增大。正切 $\tan z=\sin z/\cos z$ 只在余弦非零处按该式定义；讨论求导与积分时必须保留这一限制。

\section{复对数、主值与全纯分支}

方程 $e^w=z\ne0$ 的解不是唯一数，而是 $\log|z|+i\arg z$ 所表示的全部值。若要作微分或积分，必须选出一个在指定开集上单值定义的函数。

\begin{definition}[对数分支]
设 $U\subset\mathbb{C}\setminus\{0\}$ 为开集。满足 $e^{L(z)}=z$ 的全纯函数 $L:U\to\mathbb{C}$ 称为复对数在 $U$ 上的\keyterm{全纯分支（holomorphic branch）}。
\end{definition}

在裂平面 $U_0=\mathbb{C}\setminus(-\infty,0]$ 上取连续辐角 $-\pi<\theta<\pi$，定义\keyterm{对数主支（principal branch）}
\begin{equation}
 \operatorname{Log}z=\log|z|+i\operatorname{Arg}z,\qquad z\in U_0.
 \label{eq:ca-principal-log}
\end{equation}
这里没有包含负实轴。第\ref{chap:ca-plane}章在负实轴上给出的点态主辐角 $\pi$，不能使本式跨过该轴连续。

\begin{proposition}
\label{prop:ca-log-derivative}
$\operatorname{Log}$ 在 $U_0$ 全纯且 $(\operatorname{Log}z)'=1/z$。一般地，每个对数全纯分支都有导数 $1/z$；在同一区域上的两个分支相差一个固定的 $2\pi i$ 整数倍。
\end{proposition}
\begin{proof}
令 $r=\sqrt{x^2+y^2}$、$u=\log r$、$v=\theta$。由 $x=r\cos\theta$、$y=r\sin\theta$ 在 $r>0$ 时的局部可逆性，求导得
\[
 u_x=\frac{x}{r^2},\quad u_y=\frac{y}{r^2},\quad
 v_x=-\frac{y}{r^2},\quad v_y=\frac{x}{r^2}.
\]
在 $U_0$ 内这些偏导数连续，柯西--黎曼方程给出导数 $(x-iy)/r^2=1/z$。一般分支满足 $e^LL'=1$，故也有 $L'=1/z$。两个分支之差的导数为零，在区域上为常数；由其指数为 $1$，该常数属于 $2\pi i\mathbb{Z}$。
\end{proof}

任意非零点附近都能选择某条不穿过该点的割线，并相应选取连续辐角，所以总有局部对数。分支的存在首先是定义域问题，而不是对公式增添一个任意常数的问题。

\section{幂函数、根式与反函数}

在具有对数分支 $L$ 的区域 $U$ 上，对固定 $\alpha\in\mathbb{C}$ 定义该分支下的幂函数
\begin{equation}
 z^\alpha=\exp(\alpha L(z)),\qquad
 (z^\alpha)'=\frac{\alpha}{z}\exp(\alpha L(z)).
 \label{eq:ca-power-branch}
\end{equation}
更换 $L$ 为 $L+2k\pi i$，函数值乘以 $e^{2k\pi i\alpha}$。若 $\alpha$ 为整数，此因子恒为 $1$，与代数中的整数幂一致；若 $\alpha=1/n$，便得到 $n$ 个根的分支。主平方根是在 $U_0$ 上取 $\exp(\operatorname{Log}z/2)$，其辐角属于 $(-\pi/2,\pi/2)$。

幂运算公式必须检查分支。例如令 $z=w=e^{3\pi i/4}$，则 $z,w,zw=-i$ 均在 $U_0$，但
\[
 \operatorname{Log}z+\operatorname{Log}w=\frac{3\pi i}{2},\qquad
 \operatorname{Log}(zw)=-\frac{\pi i}{2}.
\]
两侧相差 $2\pi i$；对它们再乘非整数并取指数，差异一般不会消失。

反三角函数也须局部选择。例如在 $0$ 的充分小邻域内，取满足 $q(0)=1$ 的 $q(z)=\sqrt{1-z^2}$，并取 $q(z)+iz$ 附近满足 $L(1)=0$ 的对数，令
\[
 A(z)=-iL\bigl(q(z)+iz\bigr).
\]
由于 $(q+iz)(q-iz)=1$，可由指数形式直接验证 $\sin A(z)=z$，且 $A(0)=0$、$A'(z)=1/q(z)$。这是反正弦的一个局部分支；公式中的根与对数若另作选择，必须重新核对分支。

\section{分支存在的局部与整体问题}

\begin{proposition}
\label{prop:ca-no-global-log}
去心平面 $\mathbb{C}\setminus\{0\}$ 上不存在连续函数 $L$ 满足 $e^{L(z)}=z$，因而也不存在全局全纯对数。
\end{proposition}
\begin{proof}
若存在这样的 $L$，令 $\ell(t)=L(e^{it})$，$0\le t\le2\pi$。比较模得 $\operatorname{Re}\ell(t)=0$，比较辐角得 $\operatorname{Im}\ell(t)-t\in2\pi\mathbb{Z}$。左侧连续而值域离散，故为常数，于是 $\ell(2\pi)-\ell(0)=2\pi i$。但 $e^{2\pi i}=e^0$，同一函数在同一点的值必须相同，矛盾。
\end{proof}

这一障碍来自绕原点一周，不是原点附近的微分计算有问题。一般非零全纯函数 $f$ 的全纯对数是满足 $e^L=f$ 的全纯函数；它在每点附近存在，却未必在整个区域上存在。第\ref{chap:ca-argument}章将用闭路积分刻画全局存在条件。

\section{例题与习题}

\begin{example}[一个平方根分支的边界值]
在 $U_0$ 上取主平方根 $S$。对 $r>0$，当 $z$ 分别从上方、下方趋于 $-r$ 时，辐角分别趋于 $\pi$、$-\pi$，从而 $S(z)$ 分别趋于 $i\sqrt r$、$-i\sqrt r$。因此这一分支不能连续跨越负实轴；但在 $-r$ 的小圆盘内另选割线，仍能构造局部全纯平方根。
\end{example}

\begin{enumerate}
  \item\label{ex:ca-exp-solutions} 求 $e^z=-2$ 的全部解；求 $\sin z=0$ 与 $\cosh z=0$ 的全部解。
  \item\label{ex:ca-log-product} 设 $z,w,zw\in U_0$。以 $\operatorname{Arg}z+\operatorname{Arg}w$ 为条件，分类写出 $\operatorname{Log}(zw)$ 与 $\operatorname{Log}z+\operatorname{Log}w$ 的关系。
  \item\label{ex:ca-square-root-domain} 在上半平面选择辐角 $0<\theta<\pi$，写出平方根分支及其像集，并求导数。
  \item\label{ex:ca-no-square-root} 模仿命题\ref{prop:ca-no-global-log}，证明去心平面上不存在连续函数 $S$ 满足 $S(z)^2=z$。
\end{enumerate}

\chapter{复曲线积分与原函数}
\label{chap:ca-integration}

复积分首先是一种沿参数曲线的实变量积分。它不仅取决于端点，也可能记录曲线怎样绕过定义域中的孔洞。把参数化、积分估计与原函数联系起来，才能准确理解下一章的柯西定理。

\section{曲线积分的定义与参数化}

复值实函数 $g=u+iv$ 的积分定义为 $\int_a^b g(t)\,dt=\int_a^b u(t)\,dt+i\int_a^b v(t)\,dt$。设 $\gamma:[a,b]\to\mathbb{C}$ 分段 $C^1$，$f$ 在曲线像上连续，定义\keyterm{复曲线积分（complex contour integral）}
\begin{equation}
  \int_\gamma f(z)\,dz=\int_a^b f(\gamma(t))\gamma'(t)\,dt.
  \label{eq:ca-contour-integral}
\end{equation}
右侧按有限个光滑参数区间分别积分后相加。若 $\gamma=x+iy$、$f=u+iv$，则
\[
 \int_\gamma f(z)\,dz=\int_\gamma(u\,dx-v\,dy)
                  +i\int_\gamma(v\,dx+u\,dy).
\]
这是两个实线积分的组合。特别地，$dz$ 含有方向信息，与非负的弧长元 $|dz|$ 不同。

若 $\phi:[c,d]\to[a,b]$ 为递增的分段 $C^1$ 双射，实变量换元给出
\[
 \int_c^d f(\gamma(\phi(s)))\gamma'(\phi(s))\phi'(s)\,ds
 =\int_a^b f(\gamma(t))\gamma'(t)\,dt.
\]
所以保向重参数化不改变积分。反向曲线记为 $\gamma^{-}$，两条首尾相接曲线的连接记为 $\gamma_1*\gamma_2$，则
\[
 \int_{\gamma^{-}}f\,dz=-\int_\gamma f\,dz,\qquad
 \int_{\gamma_1*\gamma_2}f\,dz=\int_{\gamma_1}f\,dz+\int_{\gamma_2}f\,dz.
\]
积分对被积函数也具有复线性，这些性质都由定义直接得到。

\section{积分估计与极限交换}

曲线长度为 $L(\gamma)=\int_a^b|\gamma'(t)|\,dt$。复值积分仍满足三角不等式：若 $I=\int_a^b g(t)\,dt\ne0$，取 $|\lambda|=1$ 使 $\lambda I=|I|$，则
\[
 |I|=\int_a^b\operatorname{Re}(\lambda g(t))\,dt\le\int_a^b|g(t)|\,dt;
\]
$I=0$ 时也成立。应用于式\eqref{eq:ca-contour-integral}，得到基本估计。

\begin{proposition}[$ML$ 估计]
\label{prop:ca-ml}
若 $|f(z)|\le M$ 对曲线 $\gamma$ 上各点成立，则
\begin{equation}
 \left|\int_\gamma f(z)\,dz\right|
 \le\int_a^b|f(\gamma(t))|\,|\gamma'(t)|\,dt\le ML(\gamma).
 \label{eq:ca-ml}
\end{equation}
\end{proposition}

若连续函数 $f_n$ 在固定曲线像上一致收敛于 $f$，由
\[
 \left|\int_\gamma(f_n-f)\,dz\right|
 \le L(\gamma)\sup_{z\in\gamma([a,b])}|f_n(z)-f(z)|
\]
便可交换极限与积分。若曲线也随 $n$ 改变，则还须同时控制长度与误差，不能直接套用这一结论。

\begin{example}[圆弧上的量级]
对半径 $R>1$ 的上半圆 $\gamma_R$，由于 $|z^2+1|\ge R^2-1$，
\[
 \left|\int_{\gamma_R}\frac{dz}{z^2+1}\right|
 \le\frac{\pi R}{R^2-1}\longrightarrow0\quad(R\to\infty).
\]
对半径 $\varepsilon$ 的圆周，若 $|f(z)|\le C|z|^{-\alpha}$，则积分模不超过 $2\pi C\varepsilon^{1-\alpha}$。当 $\alpha<1$ 时它趋于零；$\alpha=1$ 时这一估计不再保证消失。
\end{example}

\section{原函数与路径无关}

若全纯函数 $F$ 在区域 $\Omega$ 上满足 $F'=f$，称 $F$ 为 $f$ 的\keyterm{原函数（primitive）}。当 $f$ 连续时，复链式法则及实积分基本定理给出
\begin{equation}
 \int_\gamma f(z)\,dz=F(\gamma(b))-F(\gamma(a)).
 \label{eq:ca-fundamental-integral}
\end{equation}
这里在每段光滑参数区间上对 $F(\gamma(t))$ 求导，再将端点差相加即可。

\begin{theorem}[原函数判别]
\label{thm:ca-primitive}
设 $f$ 在区域 $\Omega$ 上连续。下列条件等价：$f$ 有原函数；沿 $\Omega$ 内分段 $C^1$ 曲线的积分只依赖端点；沿 $\Omega$ 内每条分段 $C^1$ 闭曲线的积分为零。
\end{theorem}
\begin{proof}
原函数存在推出路径无关，由式\eqref{eq:ca-fundamental-integral} 得到。路径无关推出闭路积分为零，可与常值曲线比较；反之，把两条同端点曲线连接成一条闭路，即得其积分相同。

现在设积分路径无关，固定 $a\in\Omega$，定义 $F(z)=\int_a^z f(\zeta)\,d\zeta$，积分沿区域内任意分段 $C^1$ 路径进行。对固定 $z$，取小圆盘 $D(z,r)\subset\Omega$。当 $|h|<r$ 时，最后一段可选为从 $z$ 到 $z+h$ 的线段，故
\[
 \frac{F(z+h)-F(z)}h=\int_0^1f(z+th)\,dt.
\]
右侧与 $f(z)$ 的差不超过 $\sup_{0\le t\le1}|f(z+th)-f(z)|$，由连续性趋于零。因此 $F'=f$。
\end{proof}

\begin{lemma}[圆盘上的三角形判别]
\label{lem:ca-triangle-primitive}
若 $f$ 在圆盘 $D$ 上连续，并且沿每个闭三角形包含于 $D$ 的三角形正向边界积分为零，则 $f$ 在 $D$ 上有原函数。
\end{lemma}
\begin{proof}
固定 $a\in D$，令 $F(z)$ 为沿线段 $[a,z]$ 的积分。对充分小的 $h$，三角形 $a,z,z+h$ 包含于 $D$，其边界积分为零给出 $F(z+h)-F(z)=\int_{[z,z+h]}f$。沿用定理\ref{thm:ca-primitive} 最后一段的估计，得到 $F'=f$。
\end{proof}

\section{典型积分与绕数}

令 $\gamma(t)=a+re^{it}$，$0\le t\le2\pi$，其中 $r>0$。直接参数化可得，对整数 $m$，
\begin{equation}
 \int_\gamma(z-a)^m\,dz
 =ir^{m+1}\int_0^{2\pi}e^{i(m+1)t}\,dt
 =\begin{cases}2\pi i,&m=-1,\\0,&m\ne-1.\end{cases}
 \label{eq:ca-circle-monomial}
\end{equation}
这一计算不以柯西定理为前提；其中唯一不消失的 $m=-1$ 项将贯穿积分公式与留数理论。

对分段 $C^1$ 闭曲线 $\gamma$ 与不在曲线上的点 $a$，定义\keyterm{绕数（winding number）}
\begin{equation}
  \operatorname{Ind}(\gamma,a)=\frac{1}{2\pi i}\int_\gamma\frac{dz}{z-a}.
  \label{eq:ca-winding-number}
\end{equation}

\begin{proposition}
\label{prop:ca-winding-integer}
绕数为整数，等于沿闭曲线连续追踪的辐角总变化除以 $2\pi$。反向使绕数变号，首尾连接使绕数相加。
\end{proposition}
\begin{proof}
沿参数区间可写 $\gamma(t)-a=r(t)e^{i\theta(t)}$，其中 $r(t)>0$，$\theta$ 连续且分段 $C^1$。这样的辐角可按有限个小区间依次选择：每段的曲线像位于不含零的小圆盘，选择局部对数，并在相邻端点处加适当的 $2\pi$ 整数倍使辐角相接。于是
\[
 \frac{\gamma'(t)}{\gamma(t)-a}=\frac{r'(t)}{r(t)}+i\theta'(t).
\]
积分后，因起终点相同，实对数项相消，得到 $\int_\gamma dz/(z-a)=i(\theta(b)-\theta(a))$。端点辐角相差 $2\pi$ 的整数倍。反向与连接的结论由积分性质得到。
\end{proof}

正向圆周绕其内部点的绕数为 $1$，绕外部点为 $0$。可将内部点连续移到圆心，或将外部点连续移向无穷远：在不穿越曲线的连通分支上，式\eqref{eq:ca-winding-number} 随 $a$ 连续而取整数，故保持不变；无穷远处积分模由 $ML$ 估计趋于零。绕数计算的是有向总圈数，不能用主辐角的端点差代替。

\section{定义域中的孔洞与积分障碍}

函数 $1/z$ 在去心平面上全纯，但绕正向单位圆的积分为 $2\pi i$，故由定理\ref{thm:ca-primitive}，它在该区域上没有原函数。在裂平面 $U_0$ 上，它有原函数 $\operatorname{Log}z$，两点间积分就是主支对数的差。同一局部公式能否有全局原函数，取决于定义域允许哪些闭路。

三个条件须加以区分：函数在曲线上连续，足以定义积分；函数在曲线附近全纯，只给出局部微分性质；函数在曲线及其内部均全纯，才支持相应的柯西积分结论。单位圆附近的 $1/z$ 满足前两项，却在内部的原点无定义。

\section{例题与习题}

\begin{example}[路径不同导致积分不同]
求 $\int_\gamma\overline z\,dz$，其中 $\gamma$ 从 $0$ 到 $1+i$。
\end{example}
\begin{solve}
沿直线 $z=(1+i)t$、$0\le t\le1$，被积式为 $2t\,dt$，积分等于 $1$。若先沿实轴到 $1$，再沿竖直线到 $1+i$，则积分为
\[
 \int_0^1t\,dt+\int_0^1(1-it)i\,dt=1+i.
\]
差异来自路径，说明 $\overline z$ 在平面上没有原函数。不能仅凭存在实偏导数就使用端点公式。
\end{solve}

\begin{enumerate}
  \item\label{ex:ca-circle-conjugate} 设 $\gamma$ 为半径 $R>0$ 的正向圆 $|z|=R$，计算 $\int_\gamma\overline z\,dz$；再说明反向时结果如何变化。
  \item\label{ex:ca-log-integral} 设路径包含于右半平面，从 $1$ 到 $1+i$。计算 $\int_\gamma dz/z$，并说明为何与所选路径无关。
  \item\label{ex:ca-winding-path} 对 $\gamma(t)=a+re^{ikt}$（$0\le t\le2\pi$），其中 $r>0$、$k\in\mathbb{Z}$，求 $\operatorname{Ind}(\gamma,a)$；包括 $k=0$ 的情形。
  \item\label{ex:ca-arc-estimate} 设 $p>1$，大半圆 $\gamma_R$ 上的连续函数满足 $|f(z)|\le C/R^p$，其中 $C$ 与 $R$ 无关。证明 $\int_{\gamma_R}f\,dz\to0$，并用 $1/z$ 说明 $p=1$ 时结论可能失败。
\end{enumerate}

\chapter{柯西积分定理与积分公式}
\label{chap:ca-cauchy}

复可微性是局部条件，柯西理论却能由它导出闭路积分与整个区域内的函数值。证明依次经过三角形上的积分消失、局部原函数、路径变形和积分公式；全纯函数的高阶光滑性将成为结论，而不是预先假设。

\section{柯西--古尔萨定理}

\begin{theorem}[柯西--古尔萨定理，Cauchy--Goursat theorem]
\label{thm:ca-goursat}
设 $T$ 为闭三角形，$f$ 在包含 $T$ 的开集上全纯，则沿正向边界有 $\int_{\partial T}f(z)\,dz=0$。此处不要求 $f'$ 连续。
\end{theorem}
\begin{proof}
记 $I(T)=\int_{\partial T}f\,dz$。连接三边中点，将 $T$ 分成四个相似小三角形。内部边的积分相消，故至少一个小三角形 $T_1$ 满足 $|I(T_1)|\ge|I(T)|/4$。重复此过程，得到嵌套闭三角形 $T_n$，其直径、周长分别为 $d/2^n$、$\ell/2^n$，并满足
\[
 |I(T)|\le4^n|I(T_n)|.
\]
紧性及直径趋零说明 $\bigcap_nT_n=\{z_0\}$。由 $f$ 在 $z_0$ 复可微，写
\[
 f(z)=f(z_0)+f'(z_0)(z-z_0)+(z-z_0)\eta(z),\qquad \eta(z)\to0.
\]
前两项有多项式原函数，沿闭路积分为零。令 $\varepsilon_n=\sup_{z\in T_n}|\eta(z)|$，在 $z_0$ 处规定 $\eta(z_0)=0$，则 $\varepsilon_n\to0$。由 $ML$ 估计，
\[
 |I(T)|\le4^n\varepsilon_n\frac d{2^n}\frac\ell{2^n}
 =d\ell\varepsilon_n\longrightarrow0.
\]
所以 $I(T)=0$。
\end{proof}

这一证明只使用一个极限点处的一阶近似。若另有 $u,v\in C^1$，也可将复积分拆成两个实线积分，再由格林公式和柯西--黎曼方程证明结论；但该途径不能代替上面对导数不预设连续性的证明。

\begin{corollary}[局部原函数]
\label{cor:ca-local-primitive}
全纯函数在其定义域内的每个圆盘上都有原函数。
\end{corollary}
\begin{proof}
全纯函数连续，三角形边界积分由定理\ref{thm:ca-goursat} 为零，应用引理\ref{lem:ca-triangle-primitive} 即可。
\end{proof}

\section{从局部积分定理到区域上的结论}

\begin{theorem}[积分的同伦不变性]
\label{thm:ca-homotopy}
设 $f$ 在开集 $U$ 上全纯。若 $U$ 中两条分段 $C^1$ 曲线定端点同伦，则其积分相等；若两条分段 $C^1$ 闭曲线通过闭曲线同伦相连，结论也成立。
\end{theorem}
\begin{proof}
设同伦为 $H:[0,1]^2\to U$。其像紧，故可用包含于 $U$、且上面各有 $f$ 的原函数的小圆盘覆盖。由紧性与 $H$ 的一致连续性，可把参数正方形分成足够细的有限矩形网格，使每个小矩形的像包含于某一个这样的圆盘。

把每条网格边的两个像端点用直线段连接。相邻矩形共用同一条线段；由于圆盘凸，这条线段包含于两侧各自选定的圆盘内。每个小矩形所对应的四段折线的积分为零，因为它们在同一个原函数的定义域内。将这些等式相加，内部边相消，只留下外边界。上、下边的折线积分等于原曲线相应子弧的积分，因为每对子弧与线段在同一圆盘内且端点相同。定端点同伦的左右边为常值；闭曲线同伦的左右边相同且方向相反。两种情形都只剩两条待比较曲线的积分之差，因而其差为零。
\end{proof}

上述证明只用到局部原函数的存在，不要求对同伦参数求导。因此中间曲线只需连续。

\begin{corollary}[柯西积分定理，Cauchy's integral theorem]
\label{cor:ca-cauchy-simply-connected}
单连通区域上的全纯函数沿每条分段 $C^1$ 闭曲线的积分为零，且在该区域上有原函数。
\end{corollary}
\begin{proof}
将闭曲线收缩为常值曲线，应用定理\ref{thm:ca-homotopy}；再用定理\ref{thm:ca-primitive}。
\end{proof}

取 $f(z)=1/(z-a)$，还得到绕数在避开 $a$ 的闭曲线同伦下不变。围绕 $a$ 的圆周绕数为 $1$，常值曲线绕数为 $0$，所以该圆周不能在去掉 $a$ 的平面内收缩。这补足了圆环与去心平面不单连通的证明。

实际计算常用另一形式：若 $\gamma$ 为分段 $C^1$ 简单闭曲线，$f$ 在曲线及其内部闭包的某开邻域上全纯，则 $\int_\gamma f\,dz=0$。这里用到附录\ref{sec:ca-jordan-tool} 所列的平面拓扑结论：该曲线能在上述邻域内收缩。函数仅在曲线附近全纯不够。

有孔区域需保留内边界。例如 $f$ 在闭圆环 $r\le|z-a|\le R$ 的邻域全纯时，两个圆周都按逆时针取向，有
\[
 \int_{|z-a|=R}f(z)\,dz=\int_{|z-a|=r}f(z)\,dz.
\]
两条圆周可通过改变半径在圆环内同伦。若按区域边界的正向，则外圆逆时针、内圆顺时针，两项之和为零。这不能简化为外圆积分本身为零。

\section{柯西积分公式}

推导积分公式时，不能预先把某个连续延拓当成全纯函数。下面的引理只给出所需的积分结论。

\begin{lemma}
\label{lem:ca-one-exception}
设 $g$ 在开集 $U$ 连续，且除一点 $a\in U$ 外全纯，则 $g$ 在 $U$ 内每个圆盘上都有原函数。
\end{lemma}
\begin{proof}
只须验证引理\ref{lem:ca-triangle-primitive} 的三角形条件。若闭三角形不含 $a$，应用古尔萨定理。若含 $a$，从 $a$ 向顶点连线，将其分成至多三个以 $a$ 为顶点的非退化三角形。在每个子三角形的 $a$ 附近切去一个尺度为 $\varepsilon$ 的小三角形，剩余多边形可以分成有限个避开 $a$ 的三角形，其边界积分为零。由于 $g$ 在 $a$ 附近有界，切除后改变的边界积分由 $ML$ 估计为 $O(\varepsilon)$。令 $\varepsilon\to0$，得到原来各子三角形的边界积分为零，再相加即可。
\end{proof}

\begin{theorem}[柯西积分公式，Cauchy's integral formula]
\label{thm:ca-cauchy-formula}
设 $\gamma$ 为正向分段 $C^1$ 简单闭曲线，$f$ 在曲线及其内部闭包的某个开邻域 $U$ 上全纯。对每个内部点 $z$，有
\begin{equation}
  f(z)=\frac{1}{2\pi i}\int_\gamma\frac{f(\zeta)}{\zeta-z}\,d\zeta.
  \label{eq:ca-cauchy-formula}
\end{equation}
\end{theorem}
\begin{proof}
固定 $z$，在 $U$ 上定义
\[
 g(\zeta)=\begin{cases}
 \dfrac{f(\zeta)-f(z)}{\zeta-z},&\zeta\ne z,\\
 f'(z),&\zeta=z.
 \end{cases}
\]
$g$ 连续，且除 $z$ 外全纯。由引理\ref{lem:ca-one-exception}，它有局部原函数；应用定理\ref{thm:ca-homotopy} 的同一证明及曲线的可收缩性，得 $\int_\gamma g\,d\zeta=0$。故
\[
 \int_\gamma\frac{f(\zeta)}{\zeta-z}\,d\zeta
 =f(z)\int_\gamma\frac{d\zeta}{\zeta-z}
 =2\pi i\operatorname{Ind}(\gamma,z)f(z).
\]
正向简单闭曲线绕内部点的绕数为 $1$，见附录\ref{sec:ca-jordan-tool}，于是得到公式。圆周的相应绕数已在命题\ref{prop:ca-winding-integer} 后直接说明，因此圆盘情形不依赖该附录中的一般拓扑结果。
\end{proof}

公式说明边界上的全部函数值决定内部函数值。它与积分定理的区别是：核 $1/(\zeta-z)$ 在内部有一个特殊点，其积分不能直接判为零。若 $z$ 在曲线上，通常积分本身便不按普通曲线积分定义；此时不能使用该公式。

\section{高阶导数公式与柯西估计}

\begin{theorem}[高阶导数公式]
\label{thm:ca-higher-derivatives}
全纯函数在定义域上具有任意阶复导数，各阶导数仍全纯。在定理\ref{thm:ca-cauchy-formula} 的条件下，对内部点 $z$ 和整数 $n\ge0$，
\begin{equation}
 f^{(n)}(z)=\frac{n!}{2\pi i}\int_\gamma
                  \frac{f(\zeta)}{(\zeta-z)^{n+1}}\,d\zeta.
 \label{eq:ca-higher-cauchy}
\end{equation}
\end{theorem}
\begin{proof}
取紧含于曲线内部的小圆盘，令其中的点到曲线的距离至少为 $d>0$。对固定 $z$ 与 $|h|<d/2$，核的差商满足
\[
 \frac1h\left(\frac1{\zeta-z-h}-\frac1{\zeta-z}\right)
 =\frac1{(\zeta-z-h)(\zeta-z)}
 \longrightarrow\frac1{(\zeta-z)^2},
\]
且在 $\zeta\in\gamma$ 上一致收敛，因为所有分母均有统一正下界。固定曲线上的一致收敛允许将极限移入积分，得到一阶公式。对核的高次幂重复同一差商论证，每次都在较小圆盘上保留正距离，便得到任意阶导数公式。每一阶导数又具有下一阶导数，因此全纯。任意定义域中的一点均可先选一个小圆盘，所以局部结论覆盖整个定义域。
\end{proof}

若 $f$ 在闭圆盘 $\overline{D(a,R)}$ 的某个开邻域上全纯，令 $M_R=\max_{|z-a|=R}|f(z)|$，由式\eqref{eq:ca-higher-cauchy} 及圆周长度 $2\pi R$，得到\keyterm{柯西估计（Cauchy estimates）}
\begin{equation}
  |f^{(n)}(a)|\le\frac{n!M_R}{R^n}.
  \label{eq:ca-cauchy-estimate}
\end{equation}
因此边界上的函数界可以控制内部导数。估计中半径不可忽略：在同一函数界下，可用的圆盘越大，导数受到的限制越强。

定理\ref{thm:ca-higher-derivatives} 也补足了第\ref{chap:ca-holomorphic}章关于光滑性的预告。逐次使用 $f_x=f'$、$f_y=if'$，可知全纯函数的实部、虚部具有任意阶连续实偏导数，因而均调和。下一章将进一步证明它们在每点附近由幂级数表示。

\section{莫雷拉定理与积分判别}

\begin{theorem}[莫雷拉定理，Morera's theorem]
\label{thm:ca-morera}
若 $f$ 在开集 $U$ 上连续，且对每个闭三角形 $T\subset U$ 都有 $\int_{\partial T}f\,dz=0$，则 $f$ 在 $U$ 上全纯。
\end{theorem}
\begin{proof}
在任意圆盘 $D\subset U$ 上，引理\ref{lem:ca-triangle-primitive} 给出原函数 $F$，满足 $F'=f$。由于 $F$ 全纯，定理\ref{thm:ca-higher-derivatives} 说明 $F'$ 也全纯，所以 $f$ 在 $D$ 全纯。圆盘可覆盖 $U$，结论成立。
\end{proof}

柯西定理从全纯性推出积分消失，莫雷拉定理则用连续性与积分消失判别全纯性。前面的单点例外引理此时还说明：连续函数若除一个点外全纯，则在该点也全纯。更一般的可去奇点条件将在第\ref{chap:ca-singularities}章讨论。

\section{例题与习题}

\begin{example}[用导数公式计算积分]
沿正向圆 $|z|=2$，计算 $I=\int e^z/(z-1)^3\,dz$。
\end{example}
\begin{solve}
$e^z$ 在闭圆盘邻域全纯，$1$ 是圆内点。式\eqref{eq:ca-higher-cauchy} 中取 $n=2$，得 $I=(2\pi i/2!)(e^z)''|_{z=1}=\pi i e$。若把路径改为 $|z|=1/2$，被积函数在整个闭圆盘邻域全纯，由积分定理得积分为零；若改为 $|z|=1$，路径经过分母零点，不能按原式直接计算。
\end{solve}

\begin{enumerate}
  \item\label{ex:ca-cauchy-integrals} 沿正向圆 $|z|=2$，分别计算 $\int e^z/z^3\,dz$ 与 $\int\cos z/(z-1)^2\,dz$。
  \item\label{ex:ca-derivative-bound} 设 $f$ 在 $D(0,R)$ 上全纯且 $|f|\le M$。证明对 $|a|<R$，有 $|f'(a)|\le M/(R-|a|)$；注意没有假设 $f$ 在大圆边界上连续。
  \item\label{ex:ca-entire-bounded} 用柯西估计证明有界整函数必为常数，并说明这不适用于仅在一个固定圆盘内有界的全纯函数。
  \item\label{ex:ca-uniform-holomorphic} 设 $f_n$ 在开集 $U$ 上全纯，并在 $U$ 的每个紧子集上一致收敛于 $f$。用莫雷拉定理证明 $f$ 全纯，并说明为什么需要先检验 $f$ 连续。
\end{enumerate}


\part{全纯函数的结构与积分应用}

\chapter{幂级数与全纯函数的解析性}
\label{chap:ca-series}

本章把第\ref{chap:ca-cauchy}章的积分表达转化为局部代数表达，完成“全纯等价于解析”的证明。重点不只是求展开式，还包括收敛区域、逐项运算与函数列极限的条件。

\section{复数项级数与函数列收敛}
复习复数项级数的绝对收敛与柯西准则，定义逐点收敛、一致收敛和局部一致收敛（locally uniform convergence）。在开集上以每个紧子集上的一致收敛刻画局部一致收敛，介绍魏尔斯特拉斯判别法。用函数列说明逐点极限未必连续，更不能据此直接判断全纯性。

\section{幂级数及其收敛半径}
研究 $\sum_{n=0}^{\infty}a_n(z-a)^n$，其中 $a\in\mathbb{C}$、$a_n\in\mathbb{C}$。证明柯西--阿达马公式
\begin{equation}
  \frac{1}{R}=\limsup_{n\to\infty}|a_n|^{1/n},
  \label{eq:ca-radius}
\end{equation}
约定上极限为 $0$ 或 $+\infty$ 时分别对应 $R=+\infty$ 或 $R=0$。证明在收敛圆盘的紧子集上绝对一致收敛，说明边界点必须分别判断，不能把圆盘内部结论直接延伸到圆周上。

\section{逐项微分、积分与泰勒展开}
先证明幂级数在收敛圆盘内可逐项微分和积分，且导数级数有相同的收敛半径。再在柯西积分公式中把核展开为几何级数，通过一致收敛交换求和与积分，得到全纯函数的泰勒展开（Taylor expansion）及系数 $a_n=f^{(n)}(a)/n!$。说明展开唯一性，并由此完成全纯性与局部解析性的等价证明。

\section{收敛半径与复奇点}
通过 $1/(1-z)$、$1/(1+z^2)$ 等例子说明实轴附近看似良好的函数，其泰勒半径可能由非实奇点决定。准确表述收敛半径的含义：它是该中心的幂级数所定义的单值全纯函数能够覆盖的最大圆盘半径；“到最近奇点的距离”需要说明所研究的函数及延拓范围，不能把原先任意选定的定义域边界全部当成奇点。

\section{全纯函数列的极限与导数收敛}
证明全纯函数列在区域内局部一致收敛时，极限仍全纯，且每一阶导数都局部一致收敛到相应导数。先用莫雷拉定理或局部柯西公式证明全纯性，再在较小圆盘上利用积分估计控制导数差。说明此结论是复分析的刚性性质，不能照搬为一般实可微函数列的结论。

\section{例题与习题安排}
\begin{enumerate}
  \item 求若干幂级数的收敛半径，分别分析圆周上不同点的收敛行为。
  \item 在不同中心展开同一个有理函数，解释奇点位置如何改变有效圆盘。
  \item 对 $f_n(z)=z^n$ 比较单位圆盘内的局部一致收敛与整个圆盘上的一致收敛。
  \item 给出逐项微分失败的实函数列作为对照，准确指出全纯函数列结论所增加的假设。
\end{enumerate}

\chapter{零点、唯一性与全纯函数的局部结构}
\label{chap:ca-local}

本章从泰勒展开的首个非零项出发，研究全纯函数的零点、唯一性和局部映射性质。依赖第\ref{chap:ca-series}章，并为辐角原理和共形映射提供结构性工具。

\section{零点的阶与局部分解}
设 $f$ 在 $a$ 的邻域内全纯，$f(a)=0$，且 $f$ 在 $a$ 的任意小邻域内不恒为零。定义零点阶数为泰勒展开中首个非零系数的指标 $m\ge1$，证明在某邻域内有 $f(z)=(z-a)^m g(z)$，其中 $g$ 全纯且 $g(a)\ne0$。由此得到局部零点孤立性；对连通区域上不恒为零的函数，其所有零点均满足上述局部假设，要在下一节通过恒等定理补足证明。另讨论乘积、商和复合如何影响零点阶数。

\section{恒等定理与唯一性}
证明恒等定理（identity theorem）：区域内两个全纯函数若在一个于区域内部有聚点的集合上相等，则在整个区域上相等。先用泰勒展开在聚点附近推出局部恒等，再利用局部恒等点集的开闭性与区域连通性传递到整个区域。说明连通性与聚点位于内部这两项条件的作用，比较实解析函数与一般光滑函数。用泰勒系数的唯一性解释为什么一个很小的局部数据集可以决定整个连通区域上的函数。

\section{局部映射与逆函数}
从局部分解出发，借助非零因子的局部全纯对数构造局部坐标，描述非常数全纯函数在一点附近与整数次幂映射相似的结构。证明 $f'(a)\ne0$ 时存在局部全纯逆函数，并说明导数为零时可出现多重覆盖。由此推导开映射定理（open mapping theorem），不预先依赖尚未建立的辐角原理。

\section{最大模原理与最小模结论}
证明非常数全纯函数不能在区域内部取得模的局部最大值。对在有界区域闭包上连续、内部全纯的函数，说明最大值如何由边界控制。若函数不为零，则对 $1/f$ 使用最大模原理得到相应的最小模结论；用存在零点的例子解释为什么非零假设必要。

\section{刘维尔定理、代数基本定理与增长}
由柯西估计证明刘维尔定理（Liouville's theorem）：有界整函数为常数，再用反证法证明每个非常数复系数多项式至少有一个复根。进一步研究若整函数满足 $|f(z)|\le C(1+|z|)^m$，其中 $C>0$、$m$ 为非负整数，则 $f$ 是次数不超过 $m$ 的多项式。把增长条件与高阶导数消失明确联系起来。

\section{例题与习题安排}
\begin{enumerate}
  \item 用展开式判断零点阶数，避免只检查函数值或一阶导数。
  \item 比较在区域内部聚集与只在边界聚集的零点，构造后者不迫使函数恒为零的例子。
  \item 用最大模原理证明一个模估计，并说明等号在内部取得时的后果。
  \item 从整函数的多项式增长条件推出其精确的次数上界，核对常数与指数的作用。
\end{enumerate}

\chapter{洛朗级数与孤立奇点}
\label{chap:ca-singularities}

本章把圆盘中的泰勒展开推广到环域，建立奇点的分类语言。学习目标是从局部行为、级数主部和极限三个角度识别同一个奇点，并准确区分孤立奇点与分支点。

\section{环域上的洛朗展开}
设 $f$ 在环域 $r<|z-a|<R$ 上全纯，其中 $0\le r<R\le+\infty$，推导洛朗展开（Laurent expansion）及系数公式
\begin{equation}
  f(z)=\sum_{n=-\infty}^{\infty}c_n(z-a)^n,\qquad
  c_n=\frac{1}{2\pi i}\int_{|\zeta-a|=\rho}
       \frac{f(\zeta)}{(\zeta-a)^{n+1}}\,d\zeta,
  \quad r<\rho<R,
  \label{eq:ca-laurent}
\end{equation}
其中圆周取正向。将正幂部分与负幂部分分别视为级数，证明它们在较小闭环域上一致收敛，并证明展开唯一性与系数不依赖所选 $\rho$。

\section{孤立奇点的三种类型}
若 $f$ 在 $0<|z-a|<R$ 上全纯，则称 $a$ 是一个待考察的孤立奇点（isolated singularity）。根据洛朗主部是零、有限项非零还是无限项非零，分别定义可去奇点（removable singularity）、极点（pole）和本性奇点（essential singularity）。证明可去性等价于局部有界；极点等价于 $|f(z)|\to\infty$，并说明极点阶数与 $1/f$ 的零点阶数的关系。

\section{本性奇点附近的取值}
以 $e^{1/z}$ 为基本例子比较不同逼近方向，证明卡索拉蒂--魏尔斯特拉斯定理（Casorati--Weierstrass theorem）：本性奇点的每个去心邻域之像在 $\mathbb{C}$ 中稠密。皮卡大定理（great Picard theorem）只作为进阶阅读目标另行介绍，清楚区分稠密性与“除至多一个值外，每个值均被取到无穷多次”的更强结论，本章不把后者当成已证明工具。

\section{无穷远点与亚纯函数}
通过 $g(w)=f(1/w)$ 定义无穷远点附近的全纯性和奇点类型。定义亚纯函数（meromorphic function）为除离散极点外全纯、且无其他奇点的函数；此处极点集合在定义域内部没有聚点。证明整函数在无穷远点可去时为常数、在无穷远点为极点时为多项式，并进一步刻画黎曼球面上的亚纯函数为有理函数。

\section{展开区域与分支点的区别}
同一个有理函数可以在以同一点为中心的不同环域中有不同洛朗展开，要求先确定奇点所在圆周，再写展开式。对 $\log z$ 和 $z^{1/2}$，解释原点的分支现象与单值函数在完整去心圆盘上的孤立奇点分类不同；不能在没有定义单值全纯函数的区域内机械套用洛朗定理。

\section{例题与习题安排}
\begin{enumerate}
  \item 将 $1/((z-1)(z-2))$ 在以原点为中心的各个最大环域中展开，逐一标注收敛区域。
  \item 判断 $\sin z/z$、$1/(e^z-1)$、$e^{1/z}$ 在原点的类型及可确定的阶数。
  \item 用 $w=1/z$ 分析多项式、有理函数和指数函数在无穷远点的行为。
  \item 比较极点、本性奇点与分支点，说明各自定义要求的邻域是什么。
\end{enumerate}

\chapter{留数定理与实积分应用}
\label{chap:ca-residues}

本章将洛朗展开中一个系数转化为全局积分工具。依赖第\ref{chap:ca-cauchy}章与第\ref{chap:ca-singularities}章；每个应用都须同时完成路径选择、奇点检查和余项估计，不能只计算留数。

\section{留数的定义与计算方法}
定义留数（residue）$\operatorname{Res}(f,a)$ 为孤立奇点 $a$ 处洛朗展开中 $(z-a)^{-1}$ 的系数。推导简单极点的极限公式，以及 $f=g/h$、$h(a)=0$、$h'(a)\ne0$ 时的 $g(a)/h'(a)$ 公式；讨论高阶极点的求导公式与直接展开方法。说明留数为零并不意味着奇点可去。

\section{留数定理及绕数形式}
先对正向分段光滑简单闭曲线 $\gamma$ 证明留数定理。若 $f$ 在曲线及其内部闭包的某个开邻域上除内部有限个孤立奇点 $a_1,\ldots,a_N$ 外全纯，则
\begin{equation}
  \int_\gamma f(z)\,dz=2\pi i\sum_{j=1}^{N}\operatorname{Res}(f,a_j).
  \label{eq:ca-residue-theorem}
\end{equation}
通过挖去小圆并使用多边界积分定理证明该式，再推广到满足相应区域条件的闭链，加入绕数权重。介绍无穷远点留数，推导黎曼球面上的总留数关系，并注意无穷远点处的坐标变换带有微分因子。

\section{有理函数与三角型实积分}
对实轴上无极点、分母次数至少比分子高二次的有理函数，使用上半圆或下半圆闭合路径并估计大圆弧。对三角有理式积分使用 $z=e^{it}$，明确 $dt=dz/(iz)$ 及单位圆上是否存在奇点。每个例题先验证原实积分的收敛或连续性，再转入复积分。

\section{振荡积分与圆弧估计}
研究带因子 $e^{iaz}$ 的积分，其中 $a\in\mathbb{R}$，根据 $a$ 的符号选择闭合半平面。陈述一个具有明确衰减假设的若尔当引理（Jordan's lemma），完整估计指数因子与圆弧长度的共同作用。区分绝对收敛、条件收敛和仅在对称截断下存在的极限。

\section{割线积分、凹进路径与主值}
通过含 $z^\alpha$ 或 $\operatorname{Log}z$ 的钥匙孔路径，展示割线上下两侧的分支值如何产生不同因子。对实轴上的极点引入柯西主值（Cauchy principal value），明确其极限定义以及小半圆贡献。强调普通广义积分不存在时，主值存在也不能被写成普通积分的值；每个分支积分还需分别控制大小圆弧。

\section{例题与习题安排}
\begin{enumerate}
  \item 计算 $\int_{-\infty}^{\infty}(1+x^4)^{-1}\,dx$，完整列出半平面内极点及大圆弧估计。
  \item 对实数 $a>|b|$，利用单位圆变换计算 $\int_0^{2\pi}(a+b\cos t)^{-1}\,dt$。
  \item 对 $0<\alpha<1$，用钥匙孔路径求 $\int_0^{\infty}x^{\alpha-1}/(1+x)\,dx$，核对分支与端点收敛。
  \item 用 $\int_{-1}^{1}dx/x$ 比较普通广义积分与主值，解释凹进方向为什么改变小圆弧贡献。
\end{enumerate}

\chapter{辐角原理与零点计数}
\label{chap:ca-argument}

本章以对数导数为桥梁，把零点和极点的重数变成曲线绕数。它既是多项式根定位的工具，也是后续函数列零点稳定性和全纯对数存在性的重要基础。

\section{对数导数与辐角原理}
设 $f$ 在正向简单闭曲线 $\gamma$ 及其内部闭包的某个开邻域上亚纯，且在 $\gamma$ 上没有零点或极点。证明对数导数 $f'/f$ 在零点或极点处的留数分别等于其阶数或阶数的负值，从而得到辐角原理（argument principle）
\begin{equation}
  \frac{1}{2\pi i}\int_\gamma\frac{f'(z)}{f(z)}\,dz=N-P,
  \label{eq:ca-argument-principle}
\end{equation}
其中 $N,P$ 分别按重数计算内部零点数与极点数。把左侧解释为像曲线 $f\circ\gamma$ 绕原点的绕数，而非端点处主辐角数值的简单相减。

\section{鲁歇定理及其证明}
证明鲁歇定理（Rouché's theorem）：若 $f,g$ 在上述曲线及内部闭包的开邻域上全纯，且边界上 $|g|<|f|$，则 $f$ 与 $f+g$ 在内部有相同的零点总数。使用 $f+tg$、$0\le t\le1$ 的边界无零点同伦解释零点数为何不变，指出边界严格不等式与重数计数的重要性。

\section{根定位、参数扰动与应用}
训练选择主导项与闭合曲线，在圆盘、圆环或半平面中估计多项式及初等超越函数的零点数。讨论参数变化时，只要边界上始终没有零点，内部零点总数就保持不变。用边界出现零点的临界情形说明计数变化的机制。

\section{全纯对数与全纯根的存在}
对区域 $\Omega$ 上处处非零的全纯函数 $f$，证明存在全纯函数 $L$ 使 $e^L=f$，当且仅当对每条闭曲线 $\gamma\subset\Omega$ 都有 $\int_\gamma f'/f\,dz=0$。由此推出单连通区域上的非零全纯函数具有全纯对数和各阶全纯根。对有孔区域，区分全纯对数的零周期条件与 $n$ 次全纯根所需的绕数可被正整数 $n$ 整除的条件。

\section{零点在局部一致极限下的行为}
结合第\ref{chap:ca-series}章，证明若全纯函数列局部一致收敛到不恒为零的函数，则在极限函数某孤立零点周围、边界无零点的充分小圆盘中，最终保持零点的总重数。由此得到赫尔维茨定理（Hurwitz's theorem）的无零点版本：处处非零的全纯函数列在区域内的局部一致极限，要么处处非零，要么恒为零。明确“区域连通”与“局部一致”各自的作用。

\section{例题与习题安排}
\begin{enumerate}
  \item 用鲁歇定理确定 $z^5+3z+1$ 在单位圆盘内的零点数，再在较大圆盘上核对全部根数。
  \item 在边界条件失效的例子中追踪零点穿过边界的过程，解释为什么不能继续沿用原计数。
  \item 比较 $z$ 与 $z^2$ 在去心平面上的全纯对数和全纯平方根存在性。
  \item 构造处处非零但趋于恒零极限的函数列，说明赫尔维茨定理中的例外不能删去。
\end{enumerate}

\part{调和函数与共形几何}

\chapter{调和函数与圆盘上的边值问题}
\label{chap:ca-harmonic}

本章从全纯函数的实部与虚部出发，把复分析与拉普拉斯方程联系起来。目标是理解调和共轭的局部与整体差异，并完整处理圆盘上连续边界数据的狄利克雷问题。

\section{调和函数与全纯函数的实部}
定义调和函数（harmonic function）为开集上的二阶连续可微实值函数 $u$，满足 $\Delta u=u_{xx}+u_{yy}=0$。利用已经证明的全纯函数光滑性和柯西--黎曼方程，严格推出全纯函数的实部与虚部调和。说明一般两个调和函数拼成 $u+iv$ 未必全纯，还必须满足共轭关系。

\section{调和共轭的构造与拓扑障碍}
定义调和共轭（harmonic conjugate），通过积分微分形式 $-u_y\,dx+u_x\,dy$ 局部构造 $v$，使 $u+iv$ 全纯。证明在单连通区域中可整体构造，且在连通区域中共轭函数若存在则相差实常数。以 $u(z)=\log|z|$ 在去心平面上为例，说明调和性不能保证存在全局单值的调和共轭。

\section{平均值性质与极值原理}
由局部全纯表示或泊松公式推导圆周平均值性质，继而得到强最大值原理与最小值原理。讨论平均值性质对连续实函数的调和性判别，并给出证明所需的局部化步骤。比较这里的实值极值原理与第\ref{chap:ca-local}章的最大模原理，避免把二者视为相同陈述。

\section{泊松核与圆盘上的狄利克雷问题}
对单位圆上的连续实值边界函数 $\varphi$，推导泊松积分（Poisson integral）
\begin{equation}
  u(re^{i\theta})=\frac{1}{2\pi}\int_0^{2\pi}
  \frac{1-r^2}{1-2r\cos(\theta-t)+r^2}\,
  \varphi(e^{it})\,dt,\qquad 0\le r<1.
  \label{eq:ca-poisson-integral}
\end{equation}
证明该函数在圆盘内调和，并连续达到指定边界数据；用最大值原理证明在闭圆盘上连续这一函数类中的唯一性。这构成圆盘上的狄利克雷问题（Dirichlet problem）的完整解，不未经说明推广到任意不规则区域。

\section{正调和函数与共形变换}
介绍哈纳克不等式（Harnack inequality）在圆盘上的具体形式及其对正调和函数局部控制的意义。证明与全纯映射复合保持调和性，并说明用双全纯映射转移边值问题时，还需检查映射与边界数据能否延拓到相关边界。电势与二维稳态温度作为应用背景，物理解释与数学边值问题分开陈述。

\section{例题与习题安排}
\begin{enumerate}
  \item 为 $u(x,y)=x^2-y^2$ 构造调和共轭，并写出对应全纯函数。
  \item 判断 $\log|z|$ 在圆环与裂平面上是否具有全局调和共轭，给出积分依据。
  \item 求单位圆边界数据 $\cos(nt)$ 的调和延拓，其中 $n$ 为非负整数，并检查原点和边界行为。
  \item 用最大值原理证明圆盘边值解的唯一性与边界误差对内部误差的控制。
\end{enumerate}

\chapter{共形映射与典型区域}
\label{chap:ca-conformal}

本章把复导数的旋转伸缩意义发展为区域变换方法。依赖第\ref{chap:ca-local}章的局部逆函数和最大模原理；调和函数的应用另依赖第\ref{chap:ca-harmonic}章。必须始终区分局部保角与整个区域上的一一对应。

\section{局部保角与双全纯映射}
定义保定向的共形映射（conformal map），证明全纯函数在导数非零处保持相交曲线的有向夹角。定义双全纯映射（biholomorphic map）为双射全纯映射且其逆全纯，说明在单复变量中单射全纯函数的导数处处非零，因此其逆自动全纯。用 $e^z$ 说明处处非零的导数不保证全局单射。

\section{莫比乌斯变换与广义圆}
研究莫比乌斯变换（Möbius transformation）$T(z)=(az+b)/(cz+d)$，其中 $a,b,c,d\in\mathbb{C}$ 且 $ad-bc\ne0$。在黎曼球面上解释分母零点与无穷远点的像，证明圆与直线所构成的广义圆族保持不变。引入交比（cross ratio），用三个点的对应关系确定变换，并检查区域内侧映向哪一侧。

\section{圆盘与半平面的自同构}
构造上半平面 $\mathbb{H}=\{z:\operatorname{Im}z>0\}$ 与单位圆盘之间的双全纯映射。证明圆盘自同构具有形式 $e^{i\theta}(z-a)/(1-\overline a z)$，其中 $a\in\mathbb{D}$、$\theta\in\mathbb{R}$；其完整性证明将在下一节施瓦茨引理后完成。讨论指定内部点与旋转归一化如何消除映射的不唯一性。

\section{施瓦茨引理与施瓦茨--皮克不等式}
证明施瓦茨引理（Schwarz lemma）：若 $f:\mathbb{D}\to\mathbb{D}$ 全纯且 $f(0)=0$，则 $|f(z)|\le|z|$、$|f'(0)|\le1$，并刻画非零点或导数处取等号的情形。通过圆盘自同构的显式候选把一般两点移到原点，推导施瓦茨--皮克不等式（Schwarz--Pick inequality），再解释圆盘的自然距离为何被全纯自映射压缩。

\section{条带、扇形、裂域与组合映射}
以指数、适当分支的幂函数和莫比乌斯变换为基本组件，构造条带、扇形、半平面与圆盘之间的映射。每次组合均核对定义域、分支、单射性、像集和边界对应。茹科夫斯基变换（Joukowski map）作为进阶例子，展示同一公式在不同限制域上有不同的单射性。

\section{例题与习题安排}
\begin{enumerate}
  \item 构造将 $\mathbb{H}$ 映到 $\mathbb{D}$ 且把给定 $z_0\in\mathbb{H}$ 映到原点的映射，并求其逆。
  \item 将指定角度的扇形映到上半平面，再映到圆盘，解释幂函数的分支选择。
  \item 求满足归一化条件的圆盘自同构，使用施瓦茨引理证明唯一性。
  \item 分析一个“导数不为零所以全局一一对应”的错误论证，并给出正确的域限制。
\end{enumerate}


\part{大域理论与进阶选读}

\chapter{正规族与黎曼映射定理}
\label{chap:ca-normal}

本章从函数序列的紧性出发，证明一类区域之间的共形映射存在。依赖柯西估计、局部一致收敛、赫尔维茨定理和施瓦茨引理；初次学习时可先理解陈述与证明结构，再补齐紧性论证。

\section{紧集上的一致控制与等度连续}
定义函数族的局部一致有界性与等度连续性，利用略大圆盘上的柯西估计控制小圆盘上的导数。复习阿尔泽拉--阿斯科利定理（Arzelà--Ascoli theorem）的适用条件，并构造区域的紧集穷竭，为对角线抽取子列作准备。强调对每个函数分别有界不能代替整个函数族在紧集上的统一界。

\section{正规族与蒙泰尔定理}
介绍正规族（normal family）：任意序列均可抽取子列，在紧集上一致收敛到全纯函数或一致趋于无穷；后一情形指各紧集上函数模的一致发散。本章主要使用局部一致有界的函数族，因此抽出的极限均为有限值全纯函数。证明蒙泰尔定理（Montel's theorem）的局部有界版本，把等度连续、紧集穷竭和对角线方法逐步连接起来。

\section{单射函数列的极限}
由赫尔维茨定理证明单射全纯函数列的局部一致极限要么单射，要么为常数。解释为何不能只在逐点极限中把不等式直接传到极限，以及导数归一化如何排除常值退化。用圆盘自映射的序列展示像区域可能向边界退化的现象。

\section{黎曼映射定理的陈述与证明路线}
完整陈述黎曼映射定理（Riemann mapping theorem）：每个非空单连通真子域 $\Omega\subsetneq\mathbb{C}$ 都与 $\mathbb{D}$ 双全纯等价；对指定 $a\in\Omega$，再要求 $f(a)=0$、$f'(a)>0$，则映射唯一。证明分为构造非空的有界单射函数族、极值序列的正规收敛、排除退化、利用平方根构造证明满射，以及用施瓦茨引理证明唯一性。分别解释去掉真子域或单连通假设为何不成立。

\section{存在性、显式公式与边界行为}
区别存在双全纯映射、写出初等公式和数值求映射三件事。圆环作为不能直接应用定理的典型例子，说明孔洞是拓扑障碍。边界延拓仅作为有额外边界条件的专题：介绍若当区域（Jordan domain）的边界同胚延拓结论，并明确它不是任意单连通域都具有的结论，也不参与本章映射存在性的核心证明。

\section{例题与习题安排}
\begin{enumerate}
  \item 检验 $\{z^n:n\ge1\}$ 在单位圆盘中的局部有界性和正规性，并比较在更大区域中的表现。
  \item 用刘维尔定理说明 $\mathbb{C}$ 不可能与 $\mathbb{D}$ 双全纯等价。
  \item 给一个映射存在但没有在本章求出显式公式的单连通区域，区分存在性证明与构造公式。
  \item 对黎曼映射的极值证明逐步标注所用前置定理，检查是否存在循环依赖。
\end{enumerate}

\chapter{解析延拓与黎曼曲面初步}
\label{chap:ca-continuation}

本章重新审视第\ref{chap:ca-elementary}章的分支问题：一个局部函数怎样沿路径延续，绕行后为什么可能改变。依赖恒等定理、局部幂级数和区域同伦；黎曼曲面部分建立概念与典型模型，不预先展开完整的曲面理论。

\section{函数芽与直接解析延拓}
定义函数芽（germ）：在同一点附近定义、且在某个更小邻域内相等的全纯函数视为同一个局部对象。定义两个相交区域上的直接解析延拓（analytic continuation），用恒等定理证明在相连的重叠部分上相容延拓的唯一性。明确延拓是在扩展函数的定义范围，而不是仅给原公式代入新的变量值。

\section{沿路径延拓与多值现象}
用一串相互重叠的圆盘与相容函数芽定义沿路径的解析延拓，分别讨论给定路径上的唯一性和终点相同但路径不同的结果。以对数绕原点一周增加 $2\pi i$、平方根绕行后变号为基本计算。说明绕行现象由整个路径相对于奇点的关系决定，不能只看路径端点。

\section{单值化定理与同伦条件}
陈述单值化定理（monodromy theorem）：若某全纯函数芽可沿区域内从基点出发的每条路径解析延拓，则沿定端点同伦的路径得到相同的终点函数芽；特别地，在单连通区域内它确定一个单值全纯延拓。证明分解为局部相容性和同伦参数的有限覆盖论证，强调“可沿每条路径延拓”是必要的已知条件。

\section{施瓦茨反射与自然边界}
证明施瓦茨反射原理（Schwarz reflection principle）的实轴版本：函数在实轴区间上侧全纯、连同该区间连续且在区间上取实值时，可在局部用共轭公式延拓到下侧。补充几何变换后的适用情形，但不无条件反射任意边界。定义自然边界（natural boundary），将具有自然边界的幂级数作为选读例子，完整证明须另列所用的间隙定理或独立论证。

\section{黎曼曲面的局部坐标}
将黎曼曲面（Riemann surface）定义为带全纯相容局部坐标图的一维复流形，并说明其底层空间需要满足豪斯多夫性和第二可数性。构造对数的无限层模型和平方根的双层模型，讨论投影到复平面时的分支。平方根模型在分支点可用根本身作为局部坐标，借此区分曲面上的正常点与投影映射的分支点。

\section{例题与习题安排}
\begin{enumerate}
  \item 从 $z=1$ 附近的对数函数芽出发，比较绕原点不同圈数后返回的函数芽。
  \item 用两个相交圆盘上的幂级数描述一次解析延拓，验证重叠部分相等的依据。
  \item 对满足实值边界条件的具体函数实施反射，并检查延拓在实轴处的连续性与全纯性。
  \item 用关系 $w^2=z$ 建立平方根曲面的局部坐标，说明 $w=0$ 处为何不是曲面本身的奇点。
\end{enumerate}

\chapter{整函数、无穷乘积与亚纯函数构造}
\label{chap:ca-entire}

本章由“已知函数求零点”转向“给定零点或主部构造函数”，研究局部数据如何组成全平面上的函数。依赖局部一致收敛、奇点理论和解析延拓；本章主要证明复平面上的构造定理，避免预先使用下一章的逼近理论。

\section{无穷乘积及其收敛}
先定义数项无穷乘积的收敛与非零极限，再研究全纯函数的乘积。对 $\prod_n(1+u_n)$，给出在各紧集上 $\sum_n|u_n|$ 一致收敛等可检验的充分条件，并明确各因子零点与极限零点的关系。利用局部对数与级数估计证明收敛，说明无穷乘积的收敛不能仅由每个因子趋于 $1$ 推出。

\section{魏尔斯特拉斯乘积定理}
引入初等因子（elementary factor）
\[
  E_0(w)=1-w,\qquad
  E_p(w)=(1-w)\exp\left(w+\frac{w^2}{2}+\cdots+\frac{w^p}{p}\right),
  \quad p\ge1.
\]
证明魏尔斯特拉斯乘积定理（Weierstrass product theorem）的全平面版本：对无有限聚点的指定零点序列及其正整数重数，可以选择因子阶数构造恰具有这些零点的整函数。将原点的零点另提为 $z^m$，说明两个具有相同零点及重数的非零整函数之商是某个整函数的指数。

\section{米塔--莱夫勒定理与主部构造}
陈述米塔--莱夫勒定理（Mittag--Leffler theorem）的全平面版本：对无有限聚点的指定点集及每点给定的有限洛朗主部，存在恰具有这些主部、且没有其他极点的亚纯函数。通过从各主部减去适当泰勒多项式获得局部一致收敛，解释修正项为何不改变主部。说明解的不唯一性恰由可以任意加上整函数体现。

\section{三角函数的乘积与部分分式}
将 $\sin(\pi z)$ 的零点分布与典范乘积联系起来，将 $\pi\cot(\pi z)$ 的极点及留数与部分分式展开联系起来。具体展开需要说明对称配对、局部一致收敛及零点或极点以外的有效区域；证明函数相等时，还须确定可能残留的非零指数因子或整函数项，不能仅凭零点或主部相同就下结论。

\section{伽马函数与增长阶的入门}
对 $\operatorname{Re}z>0$ 定义伽马函数（Gamma function）
\[
  \Gamma(z)=\int_0^{\infty}t^{z-1}e^{-t}\,dt,
  \qquad t^{z-1}=\exp((z-1)\log t),\quad t>0.
\]
在紧子集上控制积分与导数，证明全纯性及递推关系，再通过递推关系作亚纯延拓，并用乘积方法研究零点与极点。选读部分对非常数整函数定义最大模 $M_f(r)=\max_{|z|=r}|f(z)|$ 及增长阶，介绍有限阶条件为何能进一步限制乘积分解中的指数因子；阿达马分解定理作为后续专题，陈述时补足有限阶假设。

\section{例题与习题安排}
\begin{enumerate}
  \item 比较未经修正的乘积与加入初等因子后的乘积，找出收敛改善来自哪一阶余项。
  \item 为一组明确给定、无有限聚点的零点及重数构造整函数，检查是否出现额外零点。
  \item 为整数点处指定的简单极点设计收敛修正，并验证每个点的主部不变。
  \item 推导伽马函数的递推关系并解释其延拓范围，区分起始积分的收敛域与延拓函数的定义域。
\end{enumerate}

\chapter{全纯逼近与综合专题}
\label{chap:ca-approximation}

本章用逼近问题收束全书：幂级数适合圆盘，若紧集有孔洞，应当选择怎样的有理函数逼近？最后通过综合问题连接积分、拓扑、几何和函数构造，并指出进一步学习的入口。

\section{从泰勒逼近到紧集上的逼近}
解释泰勒多项式在收敛圆盘内部的紧集上一致逼近原函数，但单个中心的幂级数未必覆盖所关心的整个紧集。定义紧集上的一致逼近，比较多项式与有理函数作为逼近函数类的差别。用环形紧集上的 $1/z$ 说明孔洞如何阻止多项式逼近，并以曲线积分给出严格矛盾。

\section{龙格逼近定理}
陈述龙格定理（Runge's theorem）的紧集版本：若 $K\subset\mathbb{C}$ 为紧集、$f$ 在 $K$ 的某个开邻域上全纯，且集合 $A\subset\mathbb{C}\setminus K$ 与补集的每个有界连通分支相交，则可以用有限极点全位于 $A$ 的有理函数在 $K$ 上一致逼近 $f$。特别地，若 $\mathbb{C}\setminus K$ 连通，则可用多项式逼近。说明有理函数允许在无穷远点有极点，即允许具有多项式部分。

\section{证明路线与条件边界}
把龙格定理的证明分为适当围道上的柯西表示、积分的有限和逼近，以及沿补集连通分支移动极点三步。每一步分别检查紧集与围道的正距离、误差的统一控制和极点位置。区别“在紧集邻域内全纯”与“仅在紧集内部全纯且在紧集上连续”；后一种条件对应更强的逼近问题，本章不将其与已证明版本混同。

\section{综合专题的组织}
设置三条综合主线。第一条从零点计数到函数构造，连接鲁歇定理、赫尔维茨定理和无穷乘积；第二条从边界数据到内部控制，连接柯西公式、最大模原理、泊松积分与正规族；第三条从路径绕行到大域障碍，连接原函数、全纯对数、单值化定理和有理逼近。每个专题均要求先列适用条件，再选择工具，而不是按计算公式归类。

\section{进一步学习的方向}
概述椭圆函数与复环面、黎曼曲面上的积分与微分、几何函数论、解析数论和多复变量的入口。对每个方向说明需要补充的知识，例如覆盖空间、微分形式、调和分析或代数；这些内容作为阅读路线，不将尚未证明的高级结果用于核心章节的论证。

\section{例题与习题安排}
\begin{enumerate}
  \item 证明 $1/z$ 不能在紧集 $\{z:1\le|z|\le2\}$ 上被多项式一致逼近，并指出允许极点在原点后发生的变化。
  \item 对多个互不相交的闭圆盘之并，检验龙格定理的条件，讨论分片给定的全纯函数如何同时逼近。
  \item 完成一个结合留数、辐角原理和参数扰动的零点计数专题，独立检查边界临界参数。
  \item 为“全纯函数的局部信息如何决定大域性质”写一份证明关系梳理，明确每个箭头所依赖的定理与条件。
\end{enumerate}


\appendix

\chapter{先修工具与常用结论索引}
\label{chap:ca-prerequisites}

本附录汇集正文调用的先修工具，并注明超出当前证明范围的拓扑结果。圆盘上的柯西理论有完整的局部证明；推广到任意简单闭曲线时，还需要平面拓扑保证其内部与绕数具有预期的性质。

\section{实分析工具}
正文采用以下实分析结论：欧氏空间中紧等价于有界且闭；紧集上的连续函数一致连续并取得实值极值；连续一阶偏导数保证实可微；实 $C^1$ 映射在雅可比行列式非零处有局部 $C^1$ 逆映射。它们分别用于紧集与边界的分离、积分极限交换、柯西--黎曼判别和局部极坐标。复值结论均按实部、虚部分解；无穷区间的积分不因局部一致收敛就自动允许交换极限。

\section{平面拓扑与紧性工具}
\label{sec:ca-jordan-tool}
开区域的折线连通性已由命题\ref{prop:ca-polygonal-connected} 证明。下列若当曲线定理（Jordan curve theorem）的形式作为平面拓扑工具引用，本书此处不证明。

\begin{theorem}[简单闭曲线的内外区域]
非退化简单闭曲线的补集恰有两个连通分支：一个有界，称为内部；另一个无界，称为外部。两者的边界均为该曲线。该曲线能在任意包含曲线及其内部的开集中收缩为一点。若曲线还分段 $C^1$，则其绕数在外部为零，在内部恒为 $1$ 或恒为 $-1$；内部绕数为 $1$ 的方向称为正向。
\end{theorem}

这一陈述及证明可参见 T.~Tao 的\href{https://terrytao.wordpress.com/2016/10/02/math-246a-notes-3-cauchys-theorem-and-its-consequences/}{Math 246A, Notes 3，Theorem 50 与附录}。不能仅凭示意图把任意自交闭曲线都划分为唯一的“内部”；一般闭路应使用各补集分支上的绕数。对圆周，这些性质已有直接计算，故第一篇的局部分析证明不依赖一般若当曲线定理。

\section{结论、条件与误用对照}
按“结论、关键假设、典型反例、正文位置”整理柯西定理、恒等定理、最大模原理、留数定理、鲁歇定理、蒙泰尔定理和黎曼映射定理。该索引随正文证明完成后补充，不在提纲阶段填入未经正文支持的编号。

\chapter{习题提示与解答}
\label{chap:ca-solutions}

以下给出第一篇各章习题的解答。题号通过正文标签引用，各节标题用于区分所属章节。

\section{复数与复平面}
\begin{description}
 \item[习题\ref{ex:ca-roots-sum}]
 $-4=4e^{i\pi}$，故四个根为 $\sqrt2e^{i(\pi/4+k\pi/2)}$（$0\le k<4$），即 $1+i,-1+i,-1-i,1-i$。若 $n$ 整除 $m$，各项 $\omega^{mk}$ 都为 $1$，和为 $n$；否则 $\omega^m\ne1$，等比求和给出 $(1-\omega^{mn})/(1-\omega^m)=0$。负整数 $m$ 也适用。
 \item[习题\ref{ex:ca-parallelogram}]
 分别展开 $|z\pm w|^2=|z|^2+|w|^2\pm2\operatorname{Re}(z\overline w)$，相加即得结论。这是平行四边形法则：两条对角线长度的平方和等于四条边长度的平方和。
 \item[习题\ref{ex:ca-domains}]
 单位圆盘和右半平面都是凸区域，因而连通且单连通。去心单位圆盘连通，但绕原点的圆周不能收缩，故不单连通。$\mathbb{C}\setminus\mathbb{R}$ 分为上、下半平面，不连通，因此不是单连通区域；两个连通分支各自单连通。
 \item[习题\ref{ex:ca-separation}]
 对每个 $a\in K$ 选 $r_a>0$，使 $D(a,2r_a)\subset U$。取 $D(a_j,r_{a_j})$ 的有限子覆盖，令 $\delta=\min_jr_{a_j}>0$。若 $z\in K$，选择覆盖它的第 $j$ 个圆盘；对 $w\notin U$，有 $|w-a_j|\ge2r_{a_j}$，故 $|w-z|\ge|w-a_j|-|z-a_j|>r_{a_j}\ge\delta$。取下确界得到所求。
\end{description}

\section{复可微性与全纯函数}
\begin{description}
 \item[习题\ref{ex:ca-affine}]
 任一点的差商都为 $a+b\overline h/h$。沿实轴与虚轴分别等于 $a+b$、$a-b$，故复可微当且仅当 $b=0$。此时处处可微且 $f'=a$；否则没有复可微点。
 \item[习题\ref{ex:ca-c1-criterion}]
 连续偏导数满足柯西--黎曼方程的条件为 $2x=\beta x$、$2\alpha y=-\beta y$，须对所有 $x,y$ 成立。因此 $\beta=2$、$\alpha=-1$，此时 $f=z^2$、$f'=2z$。
 \item[习题\ref{ex:ca-real-constant}]
 不必。定义域有两个连通分支，实值全纯函数在每个分支上各为常数，但两个常数可以不同。例如右分支取 $1$、左分支取 $0$，在每一点附近均为常数，因而全纯。
 \item[习题\ref{ex:ca-modulus-constant}]
 $c=0$ 时 $f\equiv0$。若 $c>0$，写 $f=u+iv$，对 $u^2+v^2=c^2$ 分别求偏导，再用柯西--黎曼方程，得到 $u u_x+v v_x=0$、$v u_x-u v_x=0$。关于 $u_x,v_x$ 的系数行列式为 $-(u^2+v^2)=-c^2\ne0$，故二者均为零，即 $f'=0$。由命题\ref{prop:ca-zero-derivative}，$f$ 为常数。
\end{description}

\section{初等函数与分支}
\begin{description}
 \item[习题\ref{ex:ca-exp-solutions}]
 $e^z=-2$ 的解为 $\log2+(2k+1)\pi i$。$\sin z=0$ 等价于 $e^{2iz}=1$，故 $z=k\pi$；$\cosh z=0$ 等价于 $e^{2z}=-1$，故 $z=(k+1/2)\pi i$，均取 $k\in\mathbb{Z}$。
 \item[习题\ref{ex:ca-log-product}]
 记 $s=\operatorname{Arg}z+\operatorname{Arg}w\in(-2\pi,2\pi)$。若 $-\pi<s<\pi$，两式相等；若 $\pi<s<2\pi$，则 $\operatorname{Log}(zw)=\operatorname{Log}z+\operatorname{Log}w-2\pi i$；若 $-2\pi<s<-\pi$，则应加上 $2\pi i$。$s=\pm\pi$ 会使 $zw$ 位于负实轴，与题设不符。
 \item[习题\ref{ex:ca-square-root-domain}]
 对 $z=re^{i\theta}$，$r>0$、$0<\theta<\pi$，取 $S(z)=\sqrt r\,e^{i\theta/2}$。其像恰为第一象限 $\{w:\operatorname{Re}w>0,\operatorname{Im}w>0\}$，并有 $S'(z)=1/(2S(z))$。像集的满性由将第一象限中的数平方得到。
 \item[习题\ref{ex:ca-no-square-root}]
 若存在连续平方根 $S$，则 $h(t)=S(e^{it})e^{-it/2}$ 满足 $h(t)^2=1$。连续函数 $h$ 在连通区间 $[0,2\pi]$ 上只能恒取 $1$ 或恒取 $-1$。于是 $S(e^{2\pi i})=-S(1)$，但左侧应等于 $S(1)$，而 $S(1)\ne0$，矛盾。
\end{description}

\section{复曲线积分与原函数}
\begin{description}
 \item[习题\ref{ex:ca-circle-conjugate}]
 圆周上 $\overline z=R^2/z$，故积分为 $R^2\int_\gamma dz/z=2\pi iR^2$；反向时为 $-2\pi iR^2$。也可直接用 $z=Re^{it}$ 参数化检验。
 \item[习题\ref{ex:ca-log-integral}]
 右半平面上的 $1/z$ 有原函数 $\operatorname{Log}z$，故积分为 $\operatorname{Log}(1+i)-\operatorname{Log}1=\tfrac12\log2+\tfrac{\pi i}{4}$，与路径无关。
 \item[习题\ref{ex:ca-winding-path}]
 差商 $\gamma'(t)/(\gamma(t)-a)=ik$，积分为 $2\pi ik$，所以绕数是 $k$。$k=0$ 时曲线恒等于 $a+r$，导数为零，绕数也是零。
 \item[习题\ref{ex:ca-arc-estimate}]
 $ML$ 估计给出积分模不超过 $\pi C R^{1-p}$，当 $p>1$ 时趋于零。对上半圆 $z=Re^{it}$、$0\le t\le\pi$，$\int dz/z=\int_0^\pi i\,dt=\pi i$，说明 $p=1$ 时仅有这一量级控制不足。
\end{description}

\section{柯西积分理论}
\begin{description}
 \item[习题\ref{ex:ca-cauchy-integrals}]
 由高阶导数公式，第一项为 $(2\pi i/2!)(e^z)''|_{z=0}=\pi i$，第二项为 $2\pi i(\cos z)'|_{z=1}=-2\pi i\sin1$。两处分母零点都严格位于路径内部。
 \item[习题\ref{ex:ca-derivative-bound}]
 对任意 $0<r<R-|a|$，闭圆盘 $\overline{D(a,r)}$ 包含于 $D(0,R)$，故 $|f'(a)|\le M/r$。令 $r\uparrow R-|a|$ 得到结论；无需把积分路径放到原定义域的边界上。
 \item[习题\ref{ex:ca-entire-bounded}]
 若全平面上 $|f|\le M$，则对任意 $a$、任意 $R>0$ 都有 $|f'(a)|\le M/R$。令 $R\to\infty$ 得 $f'(a)=0$，故 $f$ 为常数。在固定圆盘中不能令可用半径趋于无穷大；例如 $f(z)=z$ 在单位圆盘内有界而非常数。
 \item[习题\ref{ex:ca-uniform-holomorphic}]
 每点有一个闭圆盘包含于 $U$，在该圆盘上一致收敛使极限连续，故 $f$ 在 $U$ 连续。对任意闭三角形 $T\subset U$，在其边界上一致收敛允许交换极限与积分，于是 $\int_{\partial T}f\,dz=\lim_n\int_{\partial T}f_n\,dz=0$。莫雷拉定理遂给出全纯性；若不先验证连续性，便缺少该定理的一项假设。
\end{description}

\end{document}
